\documentclass[a4paper]{article}
\pdfoutput=1
\usepackage[utf8]{inputenc}
\usepackage{jheppub}
\usepackage{overpic}
\usepackage{subcaption}
\usepackage{fancybox}
\usepackage{float}
\usepackage{amsfonts}
\usepackage{amsmath}
\usepackage{amsbsy}
\usepackage{graphicx}
\usepackage{amssymb}
\usepackage{amsthm}
\usepackage{bm}
\usepackage{comment}
\usepackage{epsfig}
\usepackage{latexsym}
\usepackage{pdflscape}
\usepackage{color}
\usepackage{stmaryrd}
\usepackage{slashed}
\usepackage{tikz}
\usetikzlibrary{arrows.meta,decorations.markings}

\tikzset{
  mid arrow/.style={
    postaction={
      decorate,
      decoration={
        markings,
        mark=at position 0.5 with {\arrow{Stealth[length=6pt,width=6pt]}}
      }
    }
  }
}

\usepackage{cancel}
\usepackage{xcolor}
\usepackage{braket}
\usepackage{empheq}
\usepackage{graphicx} % Allows including images
\usepackage{tikz-3dplot}
% réglage de la vue 3D : élévation=20°, azimut=110°
\tdplotsetmaincoords{20}{110}
\usetikzlibrary{decorations.pathreplacing,calc}
\usetikzlibrary{decorations.pathmorphing}

\definecolor{blue5}{rgb}{0, 0, 0.66666}
\definecolor{red4}{rgb}{0.8, 0, 0}
\definecolor{navyblue}{rgb}{0, 0.2745, 0.67843}
\definecolor{cornflowerblue}{rgb}{0.388235, 0.6941176, 0.898039}
\definecolor{forestgreen}{rgb}{0.13, 0.55, 0.13}
\usetikzlibrary{decorations.text,decorations.pathmorphing}

%Removes Header------------------------------------
\makeatletter
\def\@fpheader{\relax}
\makeatother



\usepackage{orcidlink}

%\DeclareMathOperator{\tr}{tr}
\DeclareMathOperator{\pf}{pf}
\DeclareMathOperator{\Ai}{Ai}
\DeclareMathOperator{\vol}{vol}
\DeclareMathOperator{\sgn}{sgn}
\DeclareMathOperator{\sn}{sn}
\DeclareMathOperator{\cn}{cn}
\DeclareMathOperator{\dn}{dn}
\DeclareMathOperator{\cd}{cd}
\DeclareMathOperator{\Det}{Det}
\DeclareMathOperator{\sech}{sech}
% \DeclareMathOperator{\tr}{tr}

\newcommand{\beq}{\begin{eqnarray}}
\newcommand{\eeq}{\end{eqnarray}}
\newcommand{\bea}{\begin{eqnarray}}
\newcommand{\eea}{\end{eqnarray}}
\newcommand{\be}{\begin{equation}}
\newcommand{\ee}{\end{equation}}
\newcommand{\bq}{\begin{equation}}
\newcommand{\eq}{\end{equation}}
\newcommand{\unit}{1\!\!1}
\newcommand{\half}{\frac{1}{2}}
\newcommand{\nn}{\nonumber}
\newcommand{\no}{\nonumber}
\def\tb{\tilde{\beta}}
\newcommand{\mpl}{M_{\rm Pl}}
\def\k{\kappa}

\newcommand{\bto}{\bar{\tau}^{\textbf{I}}}
\newcommand{\btt}{\bar{\tau}^{\textbf{II}}}




%%% Other definitions


\def\sp{\;\;\;,\;\;\;}
\def\l{\lambda}
%\def\z{\zeta}
\def\lab{\label}
\def\f{H}
\def\La{\Lambda}
\def\z{\zeta}
\def\t{\tau}
\def\e{\epsilon}
\def\o{\omega}
\def\le{\left}
\def\ri{\right}
\def\y{\psi}
\def\half{\frac12}
\def\q{\theta}
\def\cO{{\cal O}}
\def\m{\mu}
\def\n{\nu}
\def\td{\tilde}
\def\d{\delta}
\def\6{\partial}
\def\de{\partial}
\def\del{\nabla}
\def\ls{\ell_s}
\def\a{\alpha}
\def\b{\beta}
\def\lab{\label}
\def\parl{\parallel}
\def\tom{\tilde{\o}}
\def\dA{\dot{A}}
\def\dh{\dot{h}}
\def\df{\dot{f}}
\def\dX{\dot{X}}
\def\bo{\bar{\o}}
\def\r{x}
\def\bo{b_0}
\def\br{\bar{r}}
\def\dx{\delta X^\parl}
\def\dxy{\delta X^2}
\def\dxz{\delta X^3}
\def\ol{\overline}
\def\GM#1{\textit{[GM: #1]}}
\def\g{\gamma}
\def\s{\sigma}
\def\tr{{\rm Tr}}
\def\eps{\epsilon}
\def\fb{f}
\def\rd{\rm d}
\newcommand{\bom}[1]{\fboxsep2mm\fbox{
           $ \displaystyle{ #1} $}}
           
           \newcommand{\rar}{\rightarrow}
\newcommand{\bb}{\mathbb}
\newcommand{\mc}{\mathcal}
\newcommand{\im}{\mathrm{Im}}
\newcommand{\re}{\mathrm{Re}}
\newcommand{\mi}{\sqrt{|\mathcal{I}_4|}}


\def\lab{\label}
\def\le{\left}
\def\ri{\right}
\def\6{\partial}
\def\eps{\epsilon}


% Color definitions
\newcommand{\red}[1]{{\color{red} #1 \color{black}}}
\newcommand{\blue}[1]{{\color{blue} #1 \color{black}}}
\newcommand{\green}[1]{{\color{green} #1 \color{black}}}
\newcommand{\yellow}[1]{{\color{yellow} #1 \color{black}}}
\newcommand{\purple}[1]{{\color{purple} #1 \color{black}}}


% Comment Abbreviations
\newcommand{\IT}[1]{{\red{IG: #1}}} %Ioannis
\newcommand{\PB}[1]{{\blue{PB: #1}}} % Panos
\newcommand{\OP}[1]{{\green{OP: #1}}} % Olga
\newcommand{\PG}[1]{{\purple{PG: #1}}} % Paul

\title{Notes on c=1/Liouville theory}

\author[a]{Panos Betzios\orcidlink{0000-0002-5350-9404},}
\author[b]{Paul Ghiringhelli,}
\author[b]{Olga Papadoulaki\orcidlink{0000-0001-5302-2930}}
\affiliation[a]{\href{https://www.ugent.be/we/physics-astronomy/en}{Department of Physics and Astronomy},
Ghent University, \\ Krijgslaan, 281-S9, 9000 Gent, Belgium}
\affiliation[b]{\href{https://www.cpht.polytechnique.fr/?q=en}{CPHT, CNRS, École polytechnique, Institut Polytechnique de Paris},  91120 Palaiseau, France}

\emailAdd{panos.betzios@ugent.be}
\emailAdd{paul.ghiringhelli@polytechnique.edu}
\emailAdd{olga.papadoulaki@polytechnique.edu}





\abstract{...}

\keywords{...}


\begin{document}
\maketitle
\flushbottom


%%%%%%%%%%%%%%%%%%%%%%%%%%%%%%%%%%%%%%%%%%%%%%%%%%%%%%%%%%
\cite{boutivas2025numericalcalculationentanglemententropy}\cite{Boutivas_2023}\cite{Boutivas_2024}\cite{ginsparg1993lectures2dgravity2d}\cite{Bombelli:1986rw}\cite{Bombelli:1986rw}

\section{Continuum theory}

\subsection{From the sigma model to the gauge fixed action.} We consider a free scalar coupled to two-dimensional Euclidean gravity, with a cosmological constant and a topological Einstein term on a smooth, compact two dimensional manifold $\mathcal{M}_h$ with genus $h$ and metric $g$.
\begin{equation}
  \mathcal{S}_{\sigma}^{(h,g,\Phi_0)}=\frac{1}{4\pi}\int_{\mathcal{M}_h} d^2 x \sqrt{g}\left(g^{\mu \nu}\partial_\mu X \partial_\nu X +\Phi_0\mathcal{R}+4 \pi \mu\right)\,.
\end{equation}
$\Phi_0$ is a constant and as we will see later is connected to the constant part of the linear dilaton background. The Gauss-Bonnet theorem states that $\frac{1}{4\pi}\int d^2x \sqrt{g}\mathcal{R}$ is a topological number known as the Euler characteristic $\chi_h=2-2h$. The partition function is defined through the path-integral over all possible metrics on $\mathcal{M}_{h}$ modulo Diffeomorphism and the boson $X$
\begin{equation}
  \mathcal{Z} = \sum _{h\geq0} (g_0)^{-\chi_h} \int \frac{[Dg^{(h)}_{\mu \nu}]}{\text{Vol}(\text{Diffeo})} [DX]e^{-\mathcal{S}_{\sigma}^{(h,g)}}\,,
\end{equation}
where we introduced the notation $g_0 = e^{\Phi_0}$. Note that we remove the $\Phi_0$ from $\mathcal{S}^{(\sigma,h,\Phi_0)}$ since we have extracted the topological dependent factor from the action. We can gauge fix the metric to the fiducial metric $g_{\mu \nu}=e^{2b\phi}\hat{g}_{\mu \nu}^{(\tau)}$. The price to pay is the introduction of b, c ghost fields. The measure is not invariant under the previous map. Also, the cosmological constant term yields a non trivial $\phi$ dependence (because of the cc term, even classically Weyl symmetry is broken). Combined together, they yield the Liouville action
\begin{equation}
  \mathcal{S}^{(\hat{g}_h)}_{L}=\frac{1}{4\pi}\int d^2x \sqrt{\hat{g}_h}\left(\hat{g}^{\mu \nu}_h \partial_\mu \phi \partial_\nu \phi +Q \hat{\mathcal{R}}\phi+4\pi\mu e^{2b\phi}\right)\,.
\end{equation}
In the case of $c=1$ matter boson $c_L=25,\,b=1,\,Q=2$ but we keep the general form of the Liouville action. We emphasize that the term in front of the Ricci scalar term comes from the anomaly of the measure. The remaining integral to perform is an integral over the conformal mode $\phi$ known as the Liouville field and an integral over the moduli space parametrized by $\tau$ introduced previously.
\begin{equation}
  M_{h}=\frac{\text{Metrics on }\mathcal{M}_h}{\text{Diffeo} \times \text{Weyl}}\,.
\end{equation}
Hence, if modular invariance is satisfied, the partition function finally reads
\begin{equation}\label{eq:finalgaugefixing}
  \begin{aligned}
    \mathcal{Z} = \sum _{h\geq0} (g_0)^{-\chi_h} \int d\tau \int [D\phi][DX][Db][Dc]e^{-(\mathcal{S}_{\sigma}+\mathcal{S}_{b,c}+\mathcal{S}_{L})}\,,
  \end{aligned}
\end{equation}
where $V_{h}$ is the volume of the moduli space of $M_h$. Let's focus on the Liouville sector for a moment. We can exploit that $\phi$ is integrated over to extract the $\mu$ dependence of the Liouville partition function. In fact, assuming the measure is invariant under $\phi \rightarrow \phi -\frac{\log(\mu)}{2b}$, one obtains
\begin{equation}\label{eq:KPZ}
  \mathcal{Z}^{(h)}_L(\mu) = \mu^{-\frac{Q\chi_h}{2b}}\mathcal{Z}^{(h)}_L(\mu = 1)\,.
\end{equation}

\subsection{Liouville action and Target space interpretation} The different terms appearing in the Liouville+matter action have a meaning in terms of Target space fields. Let's recall the general form of the generalized non linear sigma model in presence of a dilaton $\Phi(Y)$, a Tachyon condensate $T(Y)$ and non trivial metric $G_{\mu \nu}(Y)$.
\begin{equation}
  \mathcal{S} = \frac{1}{4\pi}\int d^2x \sqrt{\hat{g}}\left(T(Y)+\hat{\mathcal{R}}\Phi(Y)+g^{ab}G_{\mu \nu}\partial_a Y^\mu \partial_b Y^\nu\right)\,.
\end{equation}
We see by comparison to the Liouville + matter action that $Y = (X,\phi)$ and
\begin{equation}\label{eq:background}
  T = 4\pi \mu e^{2b \phi}\sp G_{\mu \nu} = \delta_{\mu \nu}\sp  \Phi = \Phi_0+Q\phi\,. 
\end{equation}
This is important to note that the background~\eqref{eq:background} is solution of the one-loop vanishing of the $\beta$ functions only for $\mu=0$. For $\mu\neq0$ the Liouville interaction should be interpreted as a tachyon-wall background rather than as a naive linear-dilaton solution of the free beta-function equations.

There is also an easy way that the potential plays the role of a Tachyon condensate. Indeed, the vertex operator reads
\begin{equation}
  \mathcal{V}_{T} = e^{i k.Y} := e^{ikX +2\alpha \phi}\,.
\end{equation}
The marginality condition fixes $\Delta_X + \Delta_\phi = \Delta_X + \alpha(Q-\alpha)=1$. Solving the previous condition with $\Delta_X=0$ and selecting the state satisfying the Seiberg bound one obtains $\alpha = b$. We denote the integrated vertex operator as
\begin{equation}
  V_{T} = \int d^2x \sqrt{\hat{g}} \,\mathcal{V}_T\,.
\end{equation}
The Tachyon condensate results in an insertion of the infinite sum of $(\mu \,V_T)^n$
\begin{equation}
  \sum_{n=0}^{\infty}\frac{(-\mu)^n}{n!} \int d^2x_1 \sqrt{\hat{g}} ... \int d^2x_n \sqrt{\hat{g}} \,e^{2 b\phi(x_1)}...e^{2 b\phi(x_n)}\,.
\end{equation}
Remark that we effectively started with a sigma model with one boson and landed on a theory which has now two bosons $X,\phi$. We say that from the non-decoupling of the conformal mode, a new dimension has emerged. This describes a inhomogeneous 2D space $(\phi,X)$. Such inhomogeneity emerges in the $\phi$ direction due to the Liouville potential. The string coupling depends on the value of the dilaton and thus on the position in the $\phi$ direction as the dilaton and the Liouville conformal mode are related by~\eqref{eq:background}. An observer sitting on the Liouville wall defined by the condition
\begin{equation}\label{eq:relationLiouvillewallcosmologicalconstant}
  4\pi\mu e^{2b\phi_\star} = 1 \Rightarrow \Phi_\star = \Phi_0 + Q\phi_\star = \Phi_0 -Q\frac{\log(4\pi\mu)}{2b}\,.
\end{equation}
observes the scattering of states with the string coupling $g_\star = e^{\Phi_\star}$. The amplitude of the process should be perturbatively described by a genus expansion in the string coupling at the position of the Liouville wall and this is exactly what the KPZ scaling relation gives (see~\eqref{eq:finalgaugefixing},~\eqref{eq:KPZ}).
\subsection{Genus 0 Lagrangian}
On the sphere $h=0$, we choose a representant of the conformal class $\hat{g}_h$ such that the metric is flat on the complex plane $(z,\bar{z})$ except at one point which is taken at infininity.
If one writes $ds^2 = 2 dz d\bar{z}$, then the Liouville Lagrangian density on the sphere is 
\begin{equation}
  \begin{aligned}
    L_{\rm bulk.} &= \frac{1}{4\pi}\left(\partial_z\phi \partial_{\bar{z}}\phi+Q \hat{\mathcal{R}}\phi+\mu e^{2b\phi}\right)\\
    &=\frac{1}{4\pi}\left(-\phi \partial_z\partial_{\bar{z}}\phi+Q \hat{\mathcal{R}}\phi+\mu e^{2b\phi}+\text{total derivative}\right)\,.
  \end{aligned}
\end{equation}
Choosing the representative of the conformal class which is flat on $\mathbb{C}$ trades the bulk curvature coupling for an asymptotic boundary condition at $z = \infty$. Thus the curvature term does not appear explicitly in the local Lagrangian density on the finite plane, but its effect is retained through the background-charge condition on $\phi$ at infinity.
\begin{equation}
  \phi(z,\bar{z}) \underset{z \to \infty}{=}  - 2Q \log (z\bar{z})+\mathcal{O}(1)\,.
\end{equation}
In fact, $\partial_z\partial_{\bar{z}}\log(z\bar{z}) = 2\pi \delta^2(z,\bar{z})$. \\
\subsection{BRST constraint/Wheeler De Witt equation and physical states}
\paragraph{BRST constraint and the Wheeler de Witt equation}
Let a theory of gravity constructed as a tensor product $M=\mathcal{H}_L \otimes \mathcal{H}_X $. We take the matter theory to be $c$ X bosons so that $c+c_L=26$. We consider the theory at genus 0 on the Riemann sphere. The X states are thus described by closed string states.  We consider a state that takes the following form 
\begin{equation}\label{eq:primarystates}
  \ket{\Psi} = \ket{\Psi_L} \otimes \ket{\Psi_X}\,,
\end{equation}
where both $\ket{\Psi_L},\,\ket{\Psi_X}$ are primary states. Upon acting with Virasoro generators, such states span the hilbert space of the gravitational theory. Nevertheless, the Hilbert space is the truncation of this bigger space by the BRST constraint
\begin{equation}
  \left(L_0^{(\phi +X)}+\bar{L}_0^{(\phi+X)}-2\right)\ket{\Psi}=0\,,
\end{equation}
where $L_0^{(\phi +X)}$ is the Virasoro generator on the plane. This condition applied on the tensor product of primary states~\eqref{eq:primarystates} imposes the marginality condition $\Delta_L +\Delta_X =1$.
We define the Wheeler-De-Witt equation as the reduction of the BRST constraint to genus $0$ minisuperspace. Violations to Wheeler-De-Witt are associated to topology changing processes and we might view the BRST constraint as incorporating all of these processes. Truncating to zero mode on the cylinder $\phi(\tau,\sigma) = \phi_0(\tau),\, \sigma \in [0,2\pi)$. The lagrangian density becomes a lagrangian
\begin{equation}
  L_{\rm zero-mode} = \frac{1}{2}\left((\partial_\tau \phi_0)^2 + 4\pi \mu e^{2b \phi_0}\right)\,.
\end{equation}
Applying canonical quantization and working in the $\phi_0$ space, one immediatly finds the associated Schrodinger equation
\begin{equation}\label{eq:Schrodinger}
 \left[-\frac{1}{2}\frac{\partial^2}{\partial{\phi_0}^2}+2\pi \mu e^{2b\phi_0}\right]\Psi_L(\phi_0)=E^{(\rm cyl)}_L \Psi_L(\phi_0)\,.
\end{equation}
We know $E^{\rm plane}_X+E^{\rm plane}_L =2$ by the BRST constraint so one might identified $E_X^{\rm plane}\to2\Delta_X,\,E^{\rm plane}_L\to 2\Delta_L$ but of course the identification between $E_X,E_L$ and $\Delta_X,\Delta_L$ is not one to one, one can shift $E^{\rm plane}_L \to E^{\rm plane}_L +a,\, E^{\rm plane}_X \to E^{\rm plane}_X -a$. The Schwarzian term has also the effect to shift the Virasoro generators
\begin{equation}
  L_0^{\rm{L,cylinder}}\to L_0^{\rm{L,plane}} -\frac{c_L}{24}\,.
\end{equation}
This fixes $E^{(\rm plane)}_L-E^{(\rm cyl)}_L = \frac{c_L}{12}$: the energy on the cylinder is shifted by the casimir energy $\frac{c_L}{12}$. Now, it appears to be convenient in the context of conformal bootstrap to parametrized
\begin{equation}
\begin{aligned}
  \Delta_L &= \frac{c_L-1}{24}-P^2 = \frac{c_L-1}{24}+p^2\,,\\
  c_L &= 1+ 6Q^2\,.
\end{aligned}
\end{equation}
The Hilbert space of Liouville theory $\mathcal{H}_L$ known as the spectrum consists of diagonal operators with $p\in \mathbb{R}(+\epsilon)$(in both spacelike and timelike case with an $\epsilon$ prescription). This means from the bootstrap perspective that the OPE rules are defined consistently with respect to the spectrum that is we have to sum over primaries with their momenta in the spectrum. If we take crossing symmetry as a requirement, we can't negociate on that. Of course, this does not mean that we can not compute expectation values of operators for which their momentum is outside the spectrum. In practice, we can access such states and such quantities by analytic continuation from the bootstrap perspective.
In order to identify $p$ with the momentum of the asymptotic states far from the Liouville wall , it is convenient to write $E^{(\rm cyl)}_L=2p^2=2\Delta_L-\frac{c}{12}+\frac{1}{12}$. Hence, the Wheeler-De-Witt equation on the Liouville sector yields
\begin{equation}
  \left[-\frac{1}{2}\frac{\partial^2}{\partial{\phi_0}^2}+2\pi \mu e^{2b\phi_0}-2p^2\right]\Psi_L(\phi_0)=0\,.
\end{equation}
Introducing the loop parameter $\ell = e^{b\phi_0}$, this becomes a modified Bessel equation
\begin{equation}\label{eq:Wheeler-De-Witt-spacelike}
  \ell^2\frac{\partial^2}{\partial{\ell}^2}\Psi_L(\ell)+\ell\frac{\partial}{\partial{\ell}}\Psi_L(\ell)-\left(\frac{4\pi\mu}{b^2}\ell^2-\frac{4p^2}{b^2}\right)\Psi_L(\ell)=0\,.
\end{equation}
States that decay in the strong coupling region $\phi_0\to \infty$ correspond to Bessel-K functions(\textbf{if $b\in \mathbb{R}$}!!). So,
\begin{equation}
  \Psi^{(p)}_L(\ell) \propto K_{\frac{2ip}{b}}\left(\frac{2\sqrt{\pi\mu}}{b}\ell\right)\,.
\end{equation}
At this stage, this is common to define $E = 2p_L$ so that in the weak coupling region for $c=1 \Rightarrow b=1$, we have a superposition of incoming and outgoing plane waves $e^{\pm i\phi E}$. As it should be, this is even in $p$ (or $E$) so that $\Delta_L$ is indeed the correct physical parameter. By definition,
\begin{equation}
  K_{\nu}(z)=\frac{I_{\nu}(z)-I_{-\nu}(z)}{\sin(\pi \nu)}\,,
\end{equation}
and we have the asymptotics $I_{\nu}(z)\underset{z\to 0}{\sim} \left(\frac{z}{2}\right)^\nu$. We further define the reflection coefficient such that asymptotically in the weak coupling region
\begin{equation}
  \Psi^{(p)}_L(\ell) \underset{\ell \to 0}{\propto} e^{2ip\phi_0}+ R(p)e^{-2ip\phi_0}\,.
\end{equation}
Due to the expression of $K_\nu$ and the asymptotic expansion of $I_\nu$, we find
\begin{equation}
  R(p)=-\sin\left(\pi\frac{2ip}{b}\right)\times \left(\frac{\pi \mu}{b^2}\right)^{2ip/b}\,.
\end{equation}
We highlight that what we described here is crystal clear in the context of spacelike Liouville theory $c_L\geq 25 \Rightarrow b\in \mathbb{R}$. In the case $c_L\leq 1\Rightarrow b \in i\mathbb{R}$, one has to analytically continue $\phi_0 \to i\phi_0,\,b\to -i\phi_0$ and the modified Bessel equation becomes a simple Bessel equation
\begin{equation}\label{eq:Wheeler-De-Witt-spacelike}
  \ell^2\frac{\partial^2}{\partial{\ell}^2}\Psi_L(\ell)+\ell\frac{\partial}{\partial{\ell}}\Psi_L(\ell)+\left(\frac{4\pi\mu}{b^2}\ell^2-\frac{4p^2}{b^2}\right)\Psi_L(\ell)=0\,.
\end{equation}
In such a case, it is still an ongoing puzzle on which boundary conditions fix the reflection coefficient.

\paragraph{Third quantized description}

The BRST constraint is the on-shell condition for states in the first quantized picture. In view of constructing a string field theory, the Green functions/propagators are the inverse of the BRST constraint operator denoted $\mathcal{K}$. The free action is constructed from the BRST constraint operator as 
\begin{equation}
  \mathcal{S}_{\rm third} = \int dY W(Y) \mathcal{K} W(Y)\,.
\end{equation}  
In non-critical bosonic string theory $Y = (X^a,\phi)$ where $X^a$ are the $c$ associated bosons and $\phi$ is the non-decoupling conformal mode of the metric while in critical string theory $c=D=26$, $\phi$ disappears. In practice, extracting $\mathcal{K}$ might be very complicated (\PG{say something about that}) but in non-critical models dual to matrix models, this operator can be infered by the euclidean propagator = two loop correlation functions. For the specific case of the $c=1$ matrix model, it was obtained that
\begin{equation}
  \mathcal{S}^{\rm c=1}_{\rm third} = \int_0^{\infty}\frac{d\ell}{\ell} \int_{-\infty}^{+\infty} dx \,W(\ell,x)(-\mathcal{H}_\ell^2-\partial_x^2) \frac{\sinh(\pi \sqrt{-\mathcal{H}_{\ell}})}{\sqrt{-\mathcal{H}_{\ell}}} W(\ell,x)\,,
\end{equation}
where $\mathcal{H}_{\ell} = \left(\ell \frac{\partial}{\partial \ell}\right)^2 - 4\mu \ell^2$ is the Wheeler-De Witt operator obtained just above. This is constructed from the Euclidean propagator
\begin{equation}
  G(E,p) = \frac{E}{\sinh(\pi E)} \frac{1}{E^2 + p^2}\,.
\end{equation}
This is interesting to obsere the pole structure. There are poles corresponding the marginality condition
\begin{equation}
  \Delta_L + \Delta_X = \frac{c_L-1}{24}+4E^2 + \frac{c_X -1}{24} + 4p^2 = 1 \Rightarrow E^2 + p^2 = 0\,.
\end{equation}
$E$ is related to the Liouville momentum by $E = 2\,p_L$ and $p$ is the momemtum associated to the $X$ close-string vertex operator. The other set of poles is due to the degenerescence of Virasoro primaries for $c_L=25 \Rightarrow b =1$. Indeed,
\begin{equation}
  E \in i \mathbb{Z} \Leftrightarrow p_L \in i \frac{\mathbb{Z}}{2}\,.
\end{equation}
This is the standard condition $p_L = i\frac{rb + s b^{-1}}{2}$ with $b=1$.
This action is non local but locality can be retrieved by switching to the $\tau$-space. It is interesting to note that the non-locality can also be written as a bilocal action. We would like to write
\begin{equation}
  \mathcal{K}W(\ell,x) = \int_{0}^{\infty}\frac{d\ell'}{\ell'} \mathcal{F}(\ell,\ell')W(\ell',x)\,.
\end{equation}
We can take advantage of the fact that $\mathcal{K}$ is easily expressed in terms of the Wheeler-De-Witt operator and that we know the eigenvector of such operator. We use delta-function normalized eigenvectors with eigenvalues $E^2$
\begin{equation}
  \psi_E(\ell)=\frac{1}{\pi}\sqrt{E\sinh(\pi E)}K_{iE}(2\sqrt{\mu}\ell)\,.
\end{equation}
The operator $\mathcal{F}$ takes formally the expression
\begin{equation}
  \mathcal{F}(\ell,\ell') =\int_{0}^{\infty}\psi_E(\ell) \frac{\sinh(\pi E)}{E}\psi_E(\ell')\,, 
\end{equation}
so that
\begin{equation}
  \mathcal{S}^{\rm c=1}_{\rm third} = \int_0^{\infty}\frac{d\ell}{\ell} \int_0^{\infty}\frac{d\ell'}{\ell'} \int_{-\infty}^{+\infty} dx \,W(\ell,x)(-\mathcal{H}_\ell^2-\partial_x^2) \mathcal{F}(\ell,\ell') W(\ell',x)\,.
\end{equation}
This bilocality signals a lack of factorization and an effective description in terms wormholes. Note that interaction terms deduced from higher order loop correlators are also expected to be important and should compete with these wormholes effects. It is unclear \PG{at least to me...} that the interacting terms give higher order corrections compared to the wormholes and that there is hierarchy of these effects.
\paragraph{Reduction of the third quantized description in genus 0 minisuperspace}
In the limit of low energy, the BRST operator exacty reduces to the Wheeler-De-Witt operator. The low energy limit restricts attention to region of the $\phi$ space before the Liouville wall. Topology changing processes and wormhole effects are thus expected to become important in the high energy regime. In the weak coupling region $\phi\to\infty$, the BRST constraint reduces to the Wheeler-De-Witt operator and locality is recovered
\begin{equation}\label{eq:simpleentanglemententropy}
  \mathcal{S}_{\rm third} = \int \rm{d}\phi \rm{d}^d X \frac{1}{2}\left(\left(\partial_\phi \Psi\right)^2+ \left(\partial_X \Psi\right)^2+4\pi \mu e^{2b\phi}\Psi^2\right)\,.
\end{equation}
\PG{If we consider a finite size systen $[\phi_1,\phi_2]$, it actually becomes important whether we impose Neumann or Dirichlet boundary conditions at the endpoints! What shall we choose ?}
We shall view this as a massive field theory whose mass is expressed in terms of the string coupling $m^2 = 2\sqrt{\pi \mu}e^{b\phi} \sim g_s(\phi)$ which depends on the position in the $\phi$ space. In particular, in the strict weak coupling region, it describes a massless scalar field theory in d+1 dimension. For the specific case of $c=1$ i.e 1+1 dimension, we therefore immediatly obtain that the entanglement entropy of a cut $A=[\phi_1,\phi_2]$ in the weak coupling region is given by Cardy-Calabrese universal law
\begin{equation}
  S^{\rm c=1}_A = \frac{1}{3}\log\left(\frac{\phi_2-\phi_1}{\epsilon_{\rm UV}}\right)+c\,,
\end{equation}
where $c$ is a non universal constant and $\epsilon_{\rm UV}$ is a UV-cutoff. Since the $\phi$-position sets the energy scale of the problem, the Liouville potential can naturally be expressed as a RG flow equation for the mass
\begin{equation}
  \frac{d m(\phi)}{d\phi} = b\, m(\phi)\,.
\end{equation}
In particular in the semi-classical limit, $b\to 0$, we have $\frac{d m(\phi)}{d\phi}\ll m(\phi)$ and the EE entropy should be exactly given by a massive Klein-Gordon theory in d+1 dimension.

Outside the strict weak coupling regime, as we will demonstrate below, the natural UV cutoff found by Hartnoll and Mazenc is naturally maped to the effect of the mass term. The mass term induces a $\mathcal{O}(1)$ effect near the Liouville wall. At such a point, it is expected that high energy effects lead to important corrections. Among them, we expect through Zamolodchikov c-theorem the central charge to converge to the Liouville central charge $c_L = 1 +6Q^2 = 25$. It is clear from~\eqref{eq:simpleentanglemententropy} that such a third quantized description can't reproduce the content of Zamolodchikov c-theorem as it breaks the famous $b\to 1/b$ duality which leaves invariant the central charge.

An other interesting observation is that such theory can be mapped to a Klein-Gordon theory by the map $w = \exp\left[b\left(\phi+iX\right)\right]$ from which we obtain
\begin{equation}
  \mathcal{S}_{\rm third} = \int \rm{d}^2w \left(2 \partial_w\Psi\partial_{\bar{w}}\Psi + \frac{2\pi \mu}{b^2} \Psi^2\right)\,.
\end{equation}
Curiously, it describes a massive scalar theory in 2D for which the mass is exactly determinded by the string coupling at the position of the Liouville wall! 
\begin{equation}
  m^2 = g_s(\Phi_\star)\,.
\end{equation}
This correspondence is not anecdotic because it is well known that the problem of defining an Hilbert in canonical quantum gravity is somewhat similar to the one faced for a Klein-Gordon theory. In two dimensions, this correspondence is explicit. Unfortunately, this correspondence does not help you making progress in corrections to Cardy-Calabrese universal law. To see that one could think of discretizing the $\phi$-space and the density matrix associated to the cut $[\phi_1,\phi_2]$ correspond to a state where we trace over the boundary conditions outside the cut.
\begin{equation}
  \rho_A=Z(\beta)^{-1}\sum_{\{\Psi(\phi)\},\phi \in \bar{A}}\int_{\Psi(\phi,0)=\Psi(\phi)}^{\Psi(\phi,\beta)=\Psi(\phi)} ... \int_{\Psi(\phi_{\rm end},0)=\Psi(\phi_{\rm end})}^{\Psi(\phi_{\rm end},\beta)=\Psi(\phi_{\rm end})} D\Psi(\phi \in A, X) e^{-S_{\rm third}}\,.
\end{equation}
This density matrix corresponds to a thermal density matrix and the pure state is obtained by taking $\beta\to0$. When going from the $\phi,X$-space to the $w$-space there is a non trivial jacobian arising from such a transformation. Clearly, this jacobian will spoil the fact that
\begin{equation}
  S^{\rm{EE},\,\phi}_{[\phi_1,\phi_2]} \neq S^{\rm{EE},\,w}_{{[e^{\phi_1},e^{\phi_2}]}}\,.
\end{equation}

\paragraph{Numerical method for computing the EE}
\paragraph{A first exercise to benchmark our numerical analysis: Entanglement entropy of a massive scalar field in 1+1d.}
We consider a massive scalar field $\phi$ in 1+1d. The hamiltonian of such system is
\begin{equation}
\begin{aligned}
     H &= \frac{1}{2}\int \rm{d}x \,\pi_\phi^2 + (\partial_x \phi)^2 + m^2 \phi^2 \\
       &= \frac{1}{2}\int \rm{d}x \pi_\phi^2 - \phi \partial^2_x \phi + m^2 \phi^2\,.
\end{aligned}
\end{equation}
The second line holds if there is no boundary or if we impose Dirichlet boundary conditions at the boundary. The entanglement entropy of a quantum field theory is not well defined as the entanglement entropy between adjoining spacetime regions is UV divergent~\cite{Witten_2018}. To circumvulate such problem, we consider a discretize/lattice version of the previous model. Let $a$ be the lattice spacing in the x-direction. Note that
\begin{equation}
  \phi(x+a)+\phi(x-a)-2\phi(x) = a^2 \frac{\partial^2 \phi}{\partial x^2}+o(a^2)\,.
\end{equation}
We can thus associate to the afforementionned QFT the following lattice model.
\begin{equation}
  H = \frac{1}{2a^2}\sum_n \pi_n^2 + (2+m^2 a^2)\phi_n^2 -(\phi_n\phi_{n+1}+\phi_n \phi_{n-1})\,.
\end{equation}
This describe a chain of nearest neighbour coupled harmonic oscillators. A method first described by Srednicky~\cite{Srednicki_1993} and Bombelli and al.~\cite{Bombelli:1986rw} can be systematically applied to such lattice system. For a recent review, see for instance~\cite{Casini_2009}\cite{bellucci2026entanglemententropymassivescalar}\cite{Boutivas_2023}\cite{Boutivas_2024}. In these preliminary works, they show that the entanglement entropy of the vacuum state for a massless scalar field in flat space is proportional to the area. The method consists in writting the Hamiltonian as
\begin{equation}
  \tilde{H} = \frac{1}{2}\sum_m \pi_n^2 + \frac{1}{2}\sum_{i\,j}\phi_i K_{ij}\phi_j\,,
\end{equation}
with
\begin{equation}
  K_{ij}= \delta_{ij}\left(\tilde{m}^2 + 2\right)-\left(\delta_{i j+1}+ \delta_{i j-1}\right)\,.
\end{equation}
To simplify notations, we droped the lattice prefactor in front of the hamiltonian and we define the lattice mass $\tilde{m}^2=m^2 a^2$. The matrix $K$ is a positive definite symmetric matrix (provided $m^2>0$). The method consists in taking its square-root on a finite but large size interval (large compared to the entanglement subregion in order to prevent side effects) compared to the subregion of the entanglement entropy. Then, the Entropy is easily computed from the eigenvalues of the $C$ matrix on the spatial subregion $A$ where $C$ is defined as
\begin{equation}
  C_{|A} = \sqrt{\braket{\phi \phi}_{|A} \braket{\pi \pi}_{|A}}  = \frac{1}{2}\sqrt{\left(K^{-1/2}\right)_{|A} \left(K^{1/2}\right)_{|A}}\,.
\end{equation}
For numerical stability, it is easier to determine the eigenvalues of $C^2_{|A}$ and then to diagonalize. Denoting $\nu_k$ the eigenvalues of $C_{|A}$, the entanglement entropy of the vacuum state reads
\begin{equation}
  S = \sum_k \left(\nu_k+\frac{1}{2}\right)\log\left(\nu_k +\frac{1}{2}\right) - \left(\nu_k-\frac{1}{2}\right)\log\left(\nu_k -\frac{1}{2}\right)
\end{equation}
\paragraph{EE in minisuperspace genus $0$ approximation.}
A discretized version of the action~\eqref{eq:simpleentanglemententropy} is simply given by ($\phi_n=n a$ is the lattice coordinate)
\begin{equation}
  H = \frac{1}{2a^2}\sum_n \pi_n^2 +(2+m_n^2 a^2)\Psi_n^2 -(\Psi_n\Psi_{n+1}+\Psi_n\Psi_{n-1})\,,
\end{equation}
with $m_n^2 = 4\pi \mu e^{2 n a b}$. For the $c=1$, after dropping the lattice prefactor, this yields
\begin{equation}
  \tilde{H} = \frac{1}{2}\sum_m \pi_n^2 + \frac{1}{2}\sum_{i\,j}\phi_i K^{L}_{ij}\phi_j\,,
\end{equation}
with
\begin{equation}
  K^{L}_{ij}= \delta_{ij}\left(4\pi \tilde{\mu} e^{2 a i} + 2\right)-\left(\delta_{i j+1}+ \delta_{i j-1}\right)\,.
\end{equation}
We introduced the lattice bare cosmological constant $\tilde{\mu} = \mu a^2$. Two things are important to notice. First, this is no more translation invariant. Second, the mass aspect depends on the lattice spacing and we can't get rid of it.
\paragraph{A note regarding Zamolodchikov c-theorem and subadditivity}
Zamolodchikov c-theorem states that for any local unitary QFT in 1+1d translationnaly and rotationnaly invariant the $c$ function is increasing~\cite{Zamolodchikov:1986gt}\cite{Casini_2007}. We emphasize here that translation in $\phi$ space is explicitly broken by the presence of this exponential factor and thus we do not expect the c-theorem to apply for the function
\begin{equation}
  c(\phi_L,\phi_R) = 3(\phi_R-\phi_L)\frac{\rm{d} S}{\rm{d}(\phi_R -\phi_L)}\,.
\end{equation}
In particular, as is it clear from the numerics, such a function would be negative at some point ruining the interpretation of the c-function as capturing the number of degrees of freedom.
\paragraph{Dependence on the choice of state.}
The Von Neumann entropy we numerically compute is associated to the vacuum and pure state $\rho = \ket{0}\bra{0}$ uniquely defined to be the state annihilated by the hamiltonian with flow parameter $t=iX$ of the third quantized action. Let us detail such procedure. We introduce an orthonormal basis of the WdW using the Kontorovich-Lebelev transform
\begin{equation}
  \begin{aligned}
    \psi_E(\ell) &= \frac{1}{\pi}\sqrt{\pi \sinh(\pi E)}K_{iE}(2\sqrt{\pi \mu}\ell)\,,\\
    \int_{0}^{+\infty} &\frac{\rm{d}\ell}{\ell}\psi_E(\ell)\psi_{E'}(\ell) = \delta(E-E')\,.
  \end{aligned}
\end{equation}
We then decompose the third quantized fields in terms of the elements of the basis
\begin{equation}
  \Psi(\phi=\log(\ell),t=iX)=\int_0^{+\infty}\frac{dE}{\sqrt{2E}} \,\psi_E(\ell) \left(a_{E}e^{-i\omega(E) t }+a^{\dagger}_{E}e^{i\omega(E) t }\right)\,.
\end{equation}
The dispersion relation is given by the WdW and yields $\omega(E)^2=E^2$. Even though the action is not translation invariant in $\phi$, it remains translation invariant in $X$. Thus, the Hamiltonian
\begin{equation}
\begin{aligned}
    H &= \int_0^{+\infty} \frac{\rm{d}\ell}{\ell} \left[\frac{1}{2}\left(\partial_t \Psi\right)^2 + \frac{\ell^2}{2}\left(\partial_\ell \Psi\right)^2 + 2\pi \mu e^{2 \phi} \Psi^2\right]\\
    &=\int_0^{+\infty} \frac{\rm{d}\ell}{\ell} \left[\frac{1}{2}\left(\partial_t \Psi\right)^2 + \frac{1}{2}\Psi\left(-\mathcal{H}_\ell \right)\Psi\right]\,.
\end{aligned}
\end{equation}
is a conserved quantity. We then rewrite the Hamiltonian in terms of the oscillator modes and we normal order to obtain
\begin{equation}
  H = \int_0^{\infty} dE E a^{\dagger}_{E} a_{E}\,.
\end{equation}
The vacuum $\ket{0}$ is uniquely defined (up to a phase) to be the state annihilated by the hamiltonian equivalently by all the oscillators. Only in this case does the point function reads
\begin{equation}
  \bra{0}\Psi(\ell,t)\Psi(\ell',t)\ket{0} =\int_0^{+\infty} \frac{dE}{2E} \Psi_E(\ell) \Psi_E'(\ell') = \frac{1}{2}\left(-\mathcal{H}_\ell\right)^{-1/2} (\ell,\ell')\,.
\end{equation}
If one describes the system in angular variable $w=e^{\phi +i X}$, define $T = \frac{w+\bar{w}}{2}, Y = \frac{w - \bar{w}}{2}$ so that to rewrite the action in the form of a KG action and subsequently define the new vacuum by asking the Haniltonian with flow parameter T to annihilated such a state, you compute EE of a different state.
\paragraph{Adding a boundary to the manifold}
Consider now the possibility of adding a boundary and a non trivial interaction.
\begin{equation}
  \mathcal{S}^{\rm bound .}_L = \frac{1}{4\pi}\int d^2x \sqrt{\hat{g}_h}\left(\hat{g}^{\mu \nu}_h \partial_\mu \phi \partial_\nu \phi +Q \hat{\mathcal{R}}\phi+4\pi\mu e^{2b\phi}\right) + \frac{1}{2\pi}\int_{\partial M}ds \sqrt{\hat{h}}\hat{K}\phi + \mu_b\int_{\partial M} ds \sqrt{\hat{h}}e^{b\phi}\,.
\end{equation}
$h$ is the physical metric at the boundary and is related to the fiducial metric via $\sqrt{h}=e^{b\phi}\sqrt{\hat{h}}$ and $\hat{K}$ is the extrinsic curvature at the boundary. $\mu_b$ is the boundary cosmological constant but should also be understood as the tension/charge of the associated brane. The variational principle is well defined as long as it satisfies Dirichlet boundary conditions at the boundary. It can also satisfies Neumann boundary conditions if
\begin{equation}
  n^{\mu}\partial_\mu\phi + \hat{K} +2\pi b \mu_b =0\,.
\end{equation}

\section{c=1 MQM}
\paragraph{From the continuum formulation to a discretize model}
In two dimensions, the Einstein–Hilbert term is topological. Therefore,it is natural to replace the continuum sum over geometries by a sum over discretized random surfaces, such as triangulations, because the discretization can correctly organize the sum over topologies. This is a special feature of two dimensions. Of course, for the continuum limit to be well defined, it should not matter the way we discretize the surface. All the descritized versions should flow to the same continuum limit. This is a crucial feature of the continuum description. The argument showed in what follows should not be taken as a rigorous proof but rather as a motivation.
For instance, we describe a two dimensional discretized connected and compact surface as a simplex of equilateral triangles, where at each vertex is attached a delta function of curvature depending on the number of triangles meeting at the vertex. The conical singularity is proportional to the deficit angle $\Delta = 2\pi - n\frac{\pi}{3}$. We may assume the following correspondence
\begin{equation}
   \int Dg^{(h)}_{\mu \nu } \rightarrow \sum_{\rm \Lambda_{h}}\,,
\end{equation}
and then the integral over the matter field $X$ in an sum over all the embedding on the real line at each vertex. Finally, denoting $a$ the size of the sides of the triangles, we may assume the following dictionary
\begin{equation}
\begin{aligned}
    \int_{\mathcal{M}_h} d^2 x \sqrt{g}g^{\mu \nu}\partial_\mu X \partial_\nu X &\rightarrow \sum_{\braket{ij}} \left(X_i-X_j\right)^2\,.\\
    \int_{\mathcal{M}_h} d^2 x \sqrt{g} \mu &\rightarrow \frac{\sqrt{3}}{4}a^2 \mu V\,.
\end{aligned}
\end{equation}
where $V$ denotes the number of vertices. We call $\kappa = \exp\left(-\frac{\sqrt{3}}{4}a^2 \mu\right)$ The partition function can thus be written as 
\begin{equation}
  \mathcal{Z} = \sum _{h\geq0} (g_0)^{-\chi_h} \sum_{\rm \Lambda_{h}} \kappa^V \prod_{i=1}^{V} \int dX_i \prod_{\braket{ij}}G(|X_i-X_j|)\,.
\end{equation}
We introduced $G(|X_i-X_j|)=\exp\left[-\left(\frac{X_i-X_j}{\sqrt{4\pi}}\right)^2\right]$\,. We may define the continuum limit as the double scaled limit where $a\to0$ and the number of vertices to infinity so that to keep the physical area fixed. At the level of the partition function this is obtained by taking $a\to0$ and $\kappa$ to a critical value $\kappa_c$ where the infinite sum shows a singular behavior. Progress can be made in evaluating such a sum by identifying it with the Feynman diagram expansion of a certain matrix model as it will be described in the following. Up to irrelevant prefactors we rewrite $\kappa = \kappa_c\exp\left(-\frac{\sqrt{3}}{4}a^2 \mu\right) $.
\paragraph{The matrix model}
Here we start with the matrix quantum mechanics defined as a euclidean path-integral over $N\times N$ hermitian matrices $M$ with cubic potential $U(M) = \rm{Tr}\left[\frac{M^2}{2}-\frac{M^3}{3!}\right]$. We also suppose that the action is defined on the compact slice of euclidean time $[-T,T]$. All in all, the partition function reads
\begin{equation}
  \mathcal{Z}(\alpha,\beta) = \int_{M(-\beta/2)=M(\beta/2)} DM(\tau) e^{-\alpha\int_{-\beta/2}^{\beta/2} d\tau \rm{Tr}\left[\frac{\dot{M}^2}{2}+ U(M)\right]}\,.
\end{equation}
\paragraph{Dictionary with the discretized sigma model}
By rescaling of the matrix we can rewrite the partition function as up to an irrelevant prefactor\footnote{The Jacobian from the field rescaling is a local normalization of the matrix path-integral measure. It contributes only analytic, scheme-dependent terms to $\log (Z)$, and is discarded when extracting the double-scaled nonanalytic free energy.}
\begin{equation}
  \mathcal{Z}(\alpha,\beta) = \int_{M(-\beta/2)=M(\beta/2)} DM(\tau) e^{-N\int_{-\beta/2}^{\beta/2} d\tau \rm{Tr}\left[\frac{\dot{M}^2}{2} + \frac{M^2}{2}-\tilde{\kappa}\frac{M^3}{3!}\right]}\,,
\end{equation}
with $\tilde{\kappa}=\sqrt{\frac{N}{\alpha}}$. Sending $\beta \to \infty$, the partition function can be formally written as an infinite sum over all  connected Feynman diagrams (t'Hooft expansion)
\begin{equation}
  \lim_{\beta\to\infty}\frac{\log(Z)}{\beta} =\sum_{h=0}^{\infty}N^{2-2h}\sum_{\Lambda}\tilde{\kappa}^V \prod_{i=1}^V \int_{-\infty}^{+\infty} d\tau_i \prod_{\braket{ij}} e^{-|\tau_i-\tau_j|}\,.
\end{equation}
The first complain is that the link factor is different from the matrix model compared to the discretized gravitational model. The equality is not microscopic; it is a universality statement about the continuum limit. We may associate the following dictionary between the discretized gravitational model and the matrix model
\begin{equation}
  \begin{aligned}
    \tilde{\kappa} &= \sqrt{\frac{N}{\alpha}} = \kappa\,,\\
    \frac{1}{N} &\sim g_0\,. 
  \end{aligned}
\end{equation}
From the discretized gravitational model, the continuum limit is obtained by sending the lattice spacing $a\to 0$. Equivalently, one tunes the bare lattice fugacity $\kappa$ to its critical value $\kappa_c$, so that the number (V) of triangles diverges while the physical area is kept fixed.

From the matrix-model perspective, it would seem that, in order to define a theory with a finite string coupling, it is necessary to consider finite $N$. This naive expectation is false: in the continuum limit the string coupling $g_0\sim 1/N$ becomes infinitely multiplicatively renormalized as $a\to0$. Thus, if $N\to \infty$ is taken simulatenously with the worldsheet continuum limit, we might obtain a theory with finite string coupling. To reach finite $\kappa_c$ while taking $N\to \infty$, $\alpha\to \infty$ which means that semi-classical WKB should be exact at the double scaled limit. While from the discretized gravitational limit, $a\to0$ corresponds to the continuum limit, there is no such thing as the size of the triangles in the Feynman diagram expansion. The remaining question is thus: what is the counter part of the continuum limit from the matrix model perspective ? We will answer this question soon. 

A useful quantity to define is the bare lattice cosmological constant $\Delta$. It should be defined as the fugacity associated to the area so
\begin{equation}
  \Delta = \kappa_c^2 - \kappa^2\,.
\end{equation}
In the continuum limit, it is related to the physical cosmological constant by $\mu \sim \frac{\Delta}{a^2}<\infty$.
\paragraph{Solving the matrix model}
Also, the partition function can be interpreted as a thermal partition of inverse temperature and thus we can write
\begin{equation}
  \mathcal{Z} = \rm{Tr}_{\mathcal{H}}\left[e^{-\beta \alpha H}\right]\,.
\end{equation}  
Here $\mathcal{H}$ denotes the Hilbert space of the theory and $H$ corresponds to the haniltonian. Setting the inverse temperature to infinity, the partition function reduces to the vacuum amplitude
\begin{equation}
  \mathcal{Z} \underset{\beta\to\infty}{\sim} e^{-\beta \alpha E_{\rm vac}}\,.
\end{equation}
In particular, one finds
\begin{equation}
  \frac{\log(Z)}{\beta} \underset{\beta\to\infty}{\to} - \alpha E_{\rm vac}\,.
\end{equation}
The objective is to compute exactly such quantity.

Now, we wish to use the symmetry under conjugation by a unitary matrix to diagonalize the matrix $M$. If we consider such symmetry as a gauge symmetry, this crucially alters the boundary condition we shall impose since configurations of M related by a unitary belong to the same gauge orbit. We postpone such subtlety for latter and consider such symmetry as a global symmetry. We diagonalize $M = \Omega_M \Lambda_M \Omega_M^\dagger$, where $\Lambda_M$ is the diagonal matrix of eigenvalues and $\Omega_M$ is a unitary matrix\PG{What is the difference between using SU(N) and U(N) ? I guess you should add the identity as a generator?}. The measure is not invariant under such change of variables and one obtains
\begin{equation}
  \begin{aligned}
    DM(\tau) &= \prod_{I=1}^N D\lambda_I(\tau)D\Omega(t) \, \Delta^2(\lambda)\,,\\
    \rm{Tr}\left[\left(\frac{dM}{d\tau}\right)^2\right] &= \rm{Tr}\left[\dot{\Lambda}^2\right] + \rm{Tr}\left[\Lambda,\dot{\Omega}_M \Omega_M^{\dagger}\right]^2 \,.
  \end{aligned}
\end{equation}
Remark that $\frac{d\Omega_M}{d\tau} \Omega_M^{\dagger}$ is Lie-algebra valued. Let's suppose we use $SU(N)$ symmetry and not $U(N)$. We can thus decompose this Lie algebra quantity in terms of the generators of $SU(N)$. It has $N^2-1$ generators and the rank is $N-1$. The Lie algebra consists in skew-hermitian traceless matrices. We take as generators of the Cartan diagonal imaginary matrices, and for the rest of the generators we consider antisymmetric real matrices $i\,T_{ij}=\delta_{ij}=-i\,T_{ji}$ and and matrices such that $iS_{ij}=i\delta_{ij}=iS_{ji}$. Hence,
\begin{equation}
  \dot{\Omega}_M \Omega_M^{\dagger} = \frac{i}{\sqrt{2}}\sum_{I<J} \dot{\alpha}_{IJ} T_{IJ} + \dot{\beta}_{IJ} S_{IJ} + \sum_{I=1}^{N-1} \dot{\alpha}_I H_I\,.
\end{equation}
We thus find
\begin{equation}
  \rm{Tr}\left[\dot{M}^2\right] = \rm{Tr}\left[\dot{\Lambda}^2\right] + \sum_{I<J} (\lambda_I-\lambda_J)^2 (\dot{\alpha}_{IJ}^2 + \dot{\beta}_{IJ}^2)\,.
\end{equation}
Discarding the volume of $SU(N)$, the partition function reads
\begin{equation}
\begin{aligned}
  \mathcal{Z}&=\int_{\lambda_{I}(-\beta/2)=\lambda_{I}(\beta/2)} \prod_{I=1}^{N}d\lambda_{I}(\tau) \Delta^2(\lambda(\tau)) \\
  &\exp\Bigg[-\alpha\int_{-\beta/2}^{\beta/2} d\tau \sum_{I=1}^N \frac{\dot{\lambda}^2_I}{2}+U(\lambda_I)+\frac{1}{2}\sum_{I<J}(\lambda_I-\lambda_J)^2(\dot{\alpha}_{ij}^2 + \dot{\beta}_{ij}^2)\Bigg]\,.
\end{aligned}
\end{equation}
Importantly, there is a remnant symmetry which is the permutation of the eigenvalues. The wavefunction should thus be symmetric under such permutation so we effectively describe bosonic particles. The important observation is that the system can be mapped to a system of $N$ fermions. Remark first that the operator 
\begin{equation}
  \mathcal{O}=\sum_{I=1}^N \frac{\pi_I^2}{2} + U(\lambda_I) +\sum_{I<J}\frac{(\pi_{IJ}^2+\tilde{\pi}^2_{IJ})}{2(\lambda_I-\lambda_J)^2}
\end{equation}
is not hermitian with respect to the non-flat measure. The $\pi$'s denote the momenta. But it is hermitian with respect to the flat measure. Nevertheless, the operator can be modifed to produce a hermitian operator with respect to the non-flat measure. In fact, if we define $H = \Delta(\lambda)^{-1} \mathcal{O}\,\Delta(\lambda)$ then $H$ is hermitian with respect to the non-flat measure i.e
\begin{equation}
  \int \prod D\lambda_I \, \Delta^2(\lambda) \, \psi_1^* H \psi_2 =  \int \prod D\lambda_I \, \Delta^2(\lambda) \,\psi_2 H \psi_1^*\,.
\end{equation}
The eigenvalue problem can thus be written as 
\begin{equation}
  \left[\sum_{I} -\frac{1}{2\alpha^2 \Delta(\lambda)} \frac{\partial^2}{\partial_{\lambda_I}^2}\Delta(\lambda) + U(\lambda_I) + \sum_{I<J} \frac{(\pi_{IJ}^2+\tilde{\pi}^2_{IJ})}{2(\lambda_I-\lambda_J)^2}\right] \psi(\lambda_1,...,\lambda_N) = E\,\psi(\lambda_1,...,\lambda_N)\,.
\end{equation}
Solving the eigenvalue problem yields an orthonormal basis with respect to the non-flat measure. Setting $\Psi = \psi \Delta(\lambda)$ produces a fermionic wavefunction. We can put in one to one correspondence to a theory of $N$ free fermions. The eigenvalue problem for the fermions reads
\begin{equation}
  \left[\sum_{I} -\frac{1}{2\alpha^2} \frac{\partial^2}{\partial_{\lambda_I}^2} + U(\lambda_I) + \sum_{I<J} \frac{(\pi_{IJ}^2+\tilde{\pi}^2_{IJ})}{2(\lambda_I-\lambda_J)^2}\right] \Psi(\lambda_1,...,\lambda_N) = E\,\Psi(\lambda_1,...,\lambda_N)\,.
\end{equation}
The lowest energy configurations do not interact as the interaction is manifestly positive. The ground state is obtained by filling the N first lowest energy levels
\begin{equation}
  E_{\rm vac} = \sum_{n=1}^N \epsilon_n\,,
\end{equation}
where $\epsilon_n$ is the $n$-th eigen-energy of
\begin{equation}
  h = -\frac{1}{2\alpha^2}\frac{\partial^2}{\partial_{\lambda}^2} + U(\lambda)\,.
\end{equation}
In the limit where the inverse planck coupling $\alpha$ goes to infinity, we can resort to a semi-classical approximation using Bohr-Sommerfeld quantization condition
\begin{equation}
  \int_{\lambda_-}^{\lambda_+} d\lambda \sqrt{2\left(\epsilon_n-U(\lambda)\right)} = \frac{\pi n}{\alpha}\,.
\end{equation}
In other words, in the semi-classical limit, the number of fermions is equal to the area of the phase space filled by the fermions
\begin{equation}
  n = \int \frac{dp d\lambda \alpha}{2\pi} \,\Theta\left(\epsilon_n-\epsilon(p,\lambda)\right)\,,
\end{equation}
with $\epsilon(p,\lambda)=\frac{p^2}{2}+U(\lambda)$. In particular, denoting the level of the fermi sea as $\mu_F$, we obtain
\begin{equation}
  \int_{\lambda_-}^{\lambda_+} d\lambda \sqrt{2\left(\mu_F-U(\lambda)\right)} = \frac{\pi N}{\alpha}=\pi \kappa^2(\mu_F)\,.
\end{equation}
$\lambda_\pm$ are the turning points of the classical orbits. Clearly, if one tunes $\mu_F$ close to the quadratic local maximum of the potential, $\kappa(\mu_F)$ develops a singular behavior. This is the counter part of the continuum limit from the matrix model perspective. It is obtained by magnificing the local quadratic maximum of the potential. Let us denote $\delta \mu_F = U_{c} -  \mu_F >0$ where $U_c=0$ is the local maximum of the potential at $\lambda =2$. The double scaled limit is obtained by $N\to \infty$ and $\delta \mu_F\to 0$ keeping the conbination finite which is the renormalized $g_s$ in the continuum or the cosmological constant $\mu := N\delta \mu_F$\footnote{The renormalized string coupling at the Liouville wall is related to the cosmological constant see~\eqref{eq:relationLiouvillewallcosmologicalconstant}}. We also introduce the density of eigenvalues
\begin{equation}
  \rho (\epsilon) = \frac{1}{\alpha}\sum \delta(\epsilon - \epsilon_i)\,.
\end{equation}
from which we also obtain
\begin{equation}
  \kappa^2(\mu_F) = \int_{0}^{\mu_F} \rho(\epsilon) d\epsilon\,.
\end{equation}
Differentiating wrt to $\delta \mu_F$, we obtain
\begin{equation}
  \frac{\partial \Delta}{\partial \delta \mu_F} = \rho(\mu_F)=\int_{\lambda_-}^{\lambda_+} \frac{d\lambda}{\sqrt{2\left(\mu_F-U(\lambda)\right)}}\,.
\end{equation}
We also have by definition 
\begin{equation}
  \alpha E_{\rm vac} = \alpha^2 \int_{0}^{\mu_F}  \rm{d}\epsilon\, \epsilon\rho(\epsilon)\,.
\end{equation}
Differentiating wrt to $\Delta$, using the chain rule and the relation obtained before yields
\begin{equation}
  \frac{\partial E_{\rm vac}}{\partial \Delta} = -\alpha \mu_F = \alpha \delta \mu_F + \text{irrelevant analytic term}\,.
\end{equation}
Tuning $\delta \mu_F\to 0$ to exhibit the singular behavior of $\rho(\mu_F)$, it was found that 
\begin{equation}
  \rho(\mu_F) = -\log(\delta \mu_F)+\mathcal{O}\left(\frac{1}{\alpha^2}\right)\,.
\end{equation}
We thus obtain
\begin{equation}
  \Delta = -\delta \mu_F \log(\delta \mu_F)+\mathcal{O}\left(\frac{1}{\alpha^2}\right)\,.
\end{equation}
We thus find at leading order in the WKB expansion
\begin{equation}
\begin{aligned}
  \alpha E_{\rm vac} &= \frac{\alpha^2 \delta \mu_F^2 \log(\delta \mu_F)}{2} +o\left(\alpha^2 \delta \mu_F^2 \log(\delta \mu_F)\right) +\mathcal{O}\left(\frac{1}{\alpha^2}\right) \\
  &= \frac{\alpha^2 \Delta^2}{2\log(\Delta)} +o\left(\frac{\alpha^2\Delta^2}{\log(\Delta)}\right) +\mathcal{O}\left(\frac{1}{\alpha^2}\right)\,.
\end{aligned}
\end{equation}
The remarkable feature of the double-scaled limit is that we can actually obtain all the terms of the partition function. The first observation is that upon rescaling the eigenvalues $\lambda \to \sqrt{\alpha}\lambda$, the eigenvalue energy problem for $h$
\begin{equation}
  \left[\frac{1}{2}\frac{\partial^2}{\partial \lambda^2}+\frac{\lambda^2}{2}-\frac{1}{\sqrt{\alpha}\frac{\lambda^3}{3!}}\right]\chi = \alpha e \chi\,,
\end{equation}
where $e = U_c - \epsilon$. First, in the limit $\alpha \to \infty$ all the higher order terms in the potential are infinitely suppressed. This is the universality property of the double-scaled limit. We could have started with a quartic potential, nothing would have changed since terms beyond the quadratic would have been infinitely suppressed. In the double-scaled limit, the hamiltonian we are interested in is
\begin{equation}
  h_{\rm DS} = -\frac{1}{2}\frac{\partial^2}{\partial \lambda^2}-\frac{\lambda^2}{2}\,,
\end{equation}
which is an inverted harmonic oscillator.
\section{Conformal bootstrap of boundary Liouville theory}
\subsection{Notations}
\paragraph{Central charge} We parametrize the central charge as
\begin{equation}
    c=1+6\,Q^2\,, \quad Q=b+b^{-1}\,.
\end{equation}
We do not assume anything particular range for the central charge. We want to obtain the shift equations as a functions of the momenta and the central charge.
\paragraph{Bulk conformal dimensions} We parametrize the conformal dimension of bulk primary fields $V_{P}(z,\bar{z})$ as 
\begin{equation}
  \Delta = \frac{Q^2}{4} - P^2\,, \quad P\in i \mathbb{R}_+\,.
\end{equation}

\paragraph{Boundary conformal dimension} We work on the upper half-plane. Boundary primary fields are denoted as $\Psi_{P}(x)^{\sigma_1\,\sigma_2}$ where $P$ is the momentum and $\sigma_{1,2}$ corresponds to the boundary conditions on the two sides of the puncture (in principle the two boundary cosmological constants.) We do not assume any particular rules for $\sigma_{i}$ (like FZZT boundary conditions).

\paragraph{Normalization} We work with normalizations of the two point functions/three point functions which are symmetric. Precisely, we write
\begin{equation}\label{eq:Normalizations}
\begin{aligned}
  \braket{V_{P_1}(0)V_{P_2}(1)}&=B(P)\, (\delta(P_1-P_2)+\delta(P_1+P_2))\,,\\
  \braket{V_{P_1}(0)V_{P_2}(1)V_{P_3}(\infty)}&= C(P_1,P_2,P_3)\,,\\
  \braket{\Psi_{P_1}^{\sigma_1 \sigma_2}(0)\Psi_{P_2}^{\sigma_2 \sigma_1}(1)}&= S_{P_1}^{\sigma_1 \sigma_2}\, (\delta(P_1-P_2)+\delta(P_1+P_2))\,,\\
  \braket{\Psi_{P_1}^{\sigma_1 \sigma_2}(0)\Psi_{P_2}^{\sigma_2 \sigma_3}(1)\Psi_{P_3}^{\sigma_3 \sigma_1}(\infty)} &= C_{P_1,P_2,P_3}^{\sigma_1 \sigma_2 \sigma_3}\,,\\
\end{aligned}
\end{equation}
and also 
\begin{equation}\label{eq:Normalizationsbis}
\begin{aligned}
  \braket{V_{P}(z,\bar{z})}&=\frac{U(P|\,\sigma)}{|z-\bar{z}|^{2\Delta}}\,.\\
  \braket{V_{P}(z,\bar{z})\Psi_{P'}^{\sigma \sigma}(x)}&=\frac{R(P;P'|\,\sigma)}{|z-\bar{z}|^{2\Delta -\Delta'}|z-x|^{2\Delta'}}
\end{aligned}
\end{equation}

\subsection{Equations from the Teschner trick}
We would like to write down all the possible equations one can get by insertion of appropriate degenerate fields. There are two things that we can do: either insert a bulk degenerate field $V_{\braket{2,1}}(z,\bar{z})$ which always decouple. This would lead to second order differential equations. We can also insert $\Psi_{\braket{3,1}}^{\sigma \sigma}(x)$ at the boundary. As was pointed out in the FZZT paper, $\Psi_{\braket{2,1}}^{\sigma \sigma}(x)$ decouples only for peculiar values of $\sigma$. The reason is that introducing a boundary modifies the ward identities obeyed by the boundary fields as the stress tensor at the boundary is perturbed by a certain primary operator which corresponds to the "line operator" (see FZZT paper: \PG{Add the expression of the boundary stress tensor in my convention}). In the other hand, there is a comprehensive reason for which $\Psi_{\braket{3,1}}^{\sigma \sigma}(x)$ always decouple. In fact, $V_{\braket{2,1}}(z,\bar{z})$ always decouple and as $\Im (z)\to 0$, the primary bulk operator can be expanded in boundary operators. The expansion should be compatible with the doubling trick\footnote{The origin of the doubling trick is encoded is the property of the stress tensor that satisfies $T(\bar{z})=\bar{T}(z)$} and thus the fusion rule of $V_{\braket{2,1}}(z)$ with itself (at $\bar{z}$ ) from which we obtain
\begin{equation}
  V_{\braket{2,1}}\,\times\,V_{\braket{2,1}} = \Psi_{\braket{1,1}}^{\sigma \sigma} + \Psi_{\braket{3,1}}^{\sigma \sigma}
\end{equation}

\paragraph{Three point function of boundary operators} We first consider the possibility of inserting a the bulk degenarate field $V_{\braket{2,1}}(z,\bar{z})$ in the boundary three point function with $\Re (z) \in(0,1)$
\begin{equation}
  \begin{aligned}
    &\braket{\Psi_{P_1}^{\sigma_1 \sigma_2}(0)V_{\braket{2,1}}(z,\bar{z})\Psi_{P_2}^{\sigma_2 \sigma_3}(1)\Psi_{P_3}^{\sigma_3 \sigma_1}(\infty)}\\
    &=\frac{R(P_{\braket{2,1}};P_{\braket{1,1}}|\sigma_2)}{S_{P_{\braket{1,1}}}^{\sigma_2 \sigma_2}}\,\frac{C_{P_1,P_{\braket{1,1}},P_1}^{\sigma_1 \sigma_2 \sigma_2}\,C_{P_1,P_2,P_3}^{\sigma_2 \sigma_2 \sigma_3}}{S_{P_1}^{\sigma_1 \sigma_2}}\,\mathcal{G}^{(1)}_{\braket{1,1}}(z,\bar{z})\\
    &+\frac{R(P_{\braket{2,1}};P_{\braket{3,1}}|\sigma_2)}{S_{P_{\braket{3,1}}}^{\sigma_2 \sigma_2}}\sum_{\eta\in\{{-1,0,1}\}}\frac{C_{P_1,P_{\braket{3,1}},P_1+\eta}^{\sigma_1 \sigma_2 \sigma_2}\,C_{P_1+\eta b,P_2,P_3}^{\sigma_1 \sigma_2 \sigma_3}}{S_{P_1+\eta b }^{\sigma_1 \sigma_2}}\,\mathcal{G}^{(1,\eta)}_{\braket{3,1}}(z,\bar{z})\\
    &=\frac{R(P_{\braket{2,1}};P_{\braket{1,1}}|\sigma_2)}{S_{P_{\braket{1,1}}}^{\sigma_2 \sigma_2}}\,\frac{C_{P_{\braket{1,1}},P_{2},P_2}^{\sigma_2 \sigma_2 \sigma_3}\,C_{P_1,P_2,P_3}^{\sigma_1 \sigma_2 \sigma_3}}{S_{P_2}^{\sigma_2 \sigma_3}}\,\mathcal{G}^{(2)}_{\braket{1,1}}(z,\bar{z})\\
    &+\frac{R(P_{\braket{2,1}};P_{\braket{3,1}}|\sigma_2)}{S_{P_{\braket{3,1}}}^{\sigma_2 \sigma_2}}\sum_{\eta\in\{{-1,0,1}\}}\frac{C_{P_{\braket{1,1}},P_{2},P_2+\eta b}^{\sigma_1 \sigma_2 \sigma_2}\,C_{P_2+\eta b,P_2,P_3}^{\sigma_1 \sigma_2 \sigma_3}}{S_{P_2+\eta b }^{\sigma_2 \sigma_3}}\,\mathcal{G}^{(2,\eta)}_{\braket{3,1}}(z,\bar{z})
  \end{aligned}
\end{equation}
The two equalities correspond to the different channel and this is the statement of crossing symmetry. The functions $\mathcal{G}$ correspond to the conformal blocks. I found the expansion \PG{To be checked!}
\begin{equation}
\begin{aligned}
  \mathcal{G}_{\braket{1,1}}^{(i)}(z,\bar{z})&=(z-\bar{z})^{\Delta'_{1,1}-2\Delta_{\braket{2,1}}}\,\left(\frac{z+\bar{z}}{2}\right)^{0} + ...\\
  \mathcal{G}_{\braket{3,1}}^{(i,\eta)}(z,\bar{z})&=(z-\bar{z})^{\Delta'_{3,1}-2\Delta_{\braket{2,1}}}\,\left(\frac{z+\bar{z}}{2}\right)^{\Delta'_{i+\eta}-\Delta'_{i}-\Delta'_{\braket{3,1}}} + ...
\end{aligned}
\end{equation}
We denoted by $\Delta'$ the conformal dimensions of the boundary primaries. The leading behavior is crucial to map them to the basis of the differential equation arising from the null vector decoupling. This generalized BPZ equation will lead to a second order (\PG{PDE I guess, since they effectively have 5 point by the doubling trick...}).

Alternatively, we can insert $\Psi^{\sigma_2 \sigma_2}_{\braket{3,1}}(x)$ with $x\in(0,1)$. One obtains
\begin{equation}
\begin{aligned}
    &\braket{\Psi_{P_1}^{\sigma_1 \sigma_2}(0)\Psi^{\sigma_2\sigma_2}_{\braket{3,1}}(x)\Psi_{P_2}^{\sigma_2 \sigma_3}(1)\Psi_{P_3}^{\sigma_3 \sigma_1}(\infty)}\\
    &=\sum_{\eta\in\{-1,0,1\}}\,\frac{C_{P_1,P_{\braket{3,1}},P_1+\eta b}^{\sigma_1 \sigma_2 \sigma_2}C_{P_1+\eta b,P_2,P_3}^{\sigma_1 \sigma_2 \sigma_3 }}{S_{P_1+\eta b}^{\sigma_1 \sigma_2}}\,\mathcal{G}^{(1,\eta)}(x)\\
    &=\sum_{\eta\in\{-1,0,1\}}\,\frac{C_{P_{\braket{3,1}},P_2,P_2+\eta b}^{\sigma_2 \sigma_2 \sigma_3}C_{P_1,P_2+\eta b,P_3}^{\sigma_1 \sigma_2 \sigma_3 }}{S_{P_2+ \eta b }^{\sigma_2 \sigma_3}}\,\mathcal{G}^{(2,\eta)}(x)
\end{aligned}
\end{equation}
The functions $\mathcal{G}$ are the conformal blocks. They have the following leading expansion
\begin{equation}
\begin{aligned}
  \mathcal{G}^{(1,\eta)}(x)\underset{x\to0}{\sim} &x^{\Delta'_{1+\eta}-\Delta'_{1}-\Delta'_{\braket{3,1}}}\\
  \mathcal{G}^{(2,\eta)}(x)\underset{x\to1}{\sim} &(1-x)^{\Delta'_{2+\eta}-\Delta'_{2}-\Delta'_{\braket{3,1}}}
\end{aligned}
\end{equation}
Unlike the previous case, the decoupling of $\Psi^{\sigma_2\sigma_2}_{\braket{3,1}}$ will produce a third order ODE (This is a ODE but we gain one order...). Let's emphasize this point: In bulk Liouville theory, the structure constants are such that they appropriate encode the connection formulae of a particular hypergeometric equation and also of higher order equations (arising from all the higher order degenerate fields decoupling). Boundary Liouville theory CFT data only encode the connection formulae of this third order equation (and those coming from fusing $\Psi^{\sigma_2\sigma_2}_{\braket{3,1}}$ with itself) while the decoupling of the null vector arise for specific values of the boundary cosmological constants. 

\paragraph{Two point function: 1 boundary operator with 1 bulk operator} It seems more convenient in this case to insert a bulk degenerate field.

\begin{equation}
  \begin{aligned}
    &\braket{V_{\braket{2,1}}(z,\bar{z})V_{P}(z_1,\bar{z}_1)\Psi_P'^{\sigma \sigma }(x)}\\
    &=\sum_{\epsilon = \pm} \frac{C(P,P_{2,1},P+\frac{\epsilon b}{2})}{B(P+\frac{\epsilon b}{2})}\,R(P+\frac{\epsilon b }{2};P'|\sigma) \mathcal{G}^{(\epsilon)}(z,z_1,x)\\
    &=\frac{R(P_{\braket{2,1}};P_{\braket{3,1}}|\sigma)}{S_{P_{\braket{3,1}}}^{\sigma \sigma}} \sum_{\eta \in \{-1,0,1\}}\frac{C_{P',P_{\braket{3,1},P'+\eta b}}^{\sigma \sigma \sigma }}{S_{P'+\eta b}^{\sigma \sigma}}R(P;P'+\eta b|\sigma)\mathcal{G}_{\braket{3,1}}^{(\eta)}(z,z_1,x)\\
    &+\frac{R(P_{\braket{2,1}};P_{\braket{1,1}}|\sigma)}{S_{P_{\braket{1,1}}}^{\sigma \sigma}} \frac{C_{P',P_{\braket{1,1},P'}}^{\sigma \sigma \sigma }}{S_{P'}^{\sigma \sigma}}R(P;P'|\sigma)\mathcal{G}_{\braket{1,1}}(z,z_1,x)\,.
  \end{aligned}
\end{equation}
\PG{I need to write the conformal blocks in terms of the invariants cross ratios (there should be two invariant cross ratios)}
The first equality corresponds to fusing the two bulk fields while the second corresponds to expanding the bulk field in terms of boundary operators.

\paragraph{1 point function}
We again insert a bulk degenerate field.
\begin{equation}
  \begin{aligned}
    \braket{V_{\braket{2,1}(z,\bar{z})}V_{P}(z,\bar{z})}
    &=\sum_{\epsilon = \pm} \frac{C_{P,P_{\braket{2,1}},P+\frac{\epsilon b }{2}}}{B(P+\frac{\epsilon b}{2})} U(P+\frac{\epsilon b}{2}|\sigma) \mathcal{G}^{(\epsilon)}(z,z_1)\\
    &=\frac{R(P_{\braket{2,1}};P_{\braket{3,1}}|\sigma)}{S_{P_{\braket{3,1}}}^{\sigma \sigma}} R(P;P_{\braket{3,1}}|\sigma)\mathcal{G}_{\braket{3,1}}(z,z_1)\\
    &+\frac{R(P_{\braket{2,1}};P_{\braket{1,1}}|\sigma)}{S_{P_{\braket{1,1}}}^{\sigma \sigma}}R(P;P_{\braket{1,1}}|\sigma)\mathcal{G}_{\braket{1,1}}(z,z_1)\,.
  \end{aligned}
\end{equation}
The first channel correspond to fusing the two bulk fields together while the second one correspond to expanding the bulk field in terms of boundary operators. 

\subsection{BPZ equation}
\paragraph{Bulk degenerate field}
The decoupling of the degenerate bulk field reads 
\begin{equation}
  \braket{\left(L_{-2}^{(z)}+\frac{1}{b^2}L_{-1}^{2\,(z)}\right)V_{\braket{2,1}}(z,\bar{z})\,X} = 0\,,
\end{equation}
where $X$ is a chain of fields inserted at different points (on the boundary or in the bulk).
\paragraph{Three point function of boundary operators}
We denote
\begin{equation}
  \mathcal{V}_{4}(z,\bar{z})=\braket{\Psi_{P_1}^{\sigma_1 \sigma_2}(0)V_{\braket{2,1}}(z,\bar{z})\Psi_{P_2}^{\sigma_2 \sigma_3}(1)\Psi_{P_3}^{\sigma_3 \sigma_1}(\infty)}\,.
\end{equation}
The constraint reads
\begin{equation}\label{eq:3pointfunctionBPZ1}
  0=\braket{\Psi_{P_1}^{\sigma_1 \sigma_2}(0)\left(L_{-2}^{(z)}+\frac{1}{b^2}L_{-1}^{2\,(z)}\right)V_{\braket{2,1}}(z,\bar{z})\Psi_{P_2}^{\sigma_2 \sigma_3}(1)\Psi_{P_3}^{\sigma_3 \sigma_1}(\infty)}\,.
\end{equation}
$L_{-1}$ just yields $\partial_z$ while $L_{-2}^{(z)}=\frac{1}{2\pi i}\int_{w \sim z}\frac{dw\,T(w)}{w-z}$. The standard trick is to deform the contour to pick up the residues at the other poles. In the meantime, we want to avoid factor of $\partial_{x_i}$ arising from the fact that 
\begin{equation}
  T(w)\Psi_{P}^{\sigma \sigma}(x_i)\sim \frac{\Delta_i'}{(w-x_{i})^2}+\frac{\partial_{x_{i}}}{w-x_{i}}\,.
\end{equation}
To do that we write
\begin{equation}
  \frac{1}{w-z}=\frac{w(w-1)}{z(z-1)}\frac{1}{w-z}-\frac{w-z}{z(z-1)}-\frac{2z-1}{z(z-1)}\,.
\end{equation}
This partial fraction decomposition has the benefit of removing one order of the pole at $w\sim0,w\sim1$ regarding the first term. The second term will correspond to the dilation operator $L_0$ while the last one will yield $L_{-1}=\partial_z$. Inserting this back to~\eqref{eq:3pointfunctionBPZ1} produces
\begin{equation}
\begin{aligned}
  0&=\frac{1}{b^2}\partial_z^2\,\mathcal{V}_{4}-\frac{\Delta_{2,1}}{z(z-1)}\mathcal{V}_{4}-\frac{2z-1}{z(z-1)}\partial_z \mathcal{V}_{4}\\
  &+\braket{\Psi_{P_1}^{\sigma 1 \sigma 2}(0)\frac{1}{2\pi i}\int_{w \sim z}\frac{w(w-1)}{z(z-1)}\frac{T(w)}{w-z} V_{\braket{2,1}}(z,\bar{z})\Psi_{P_2}^{\sigma 2 \sigma 3}(1)\Psi_{P_3}^{\sigma 3 \sigma 1}(\infty)}\,.
\end{aligned}
\end{equation}
We then deform the contour to pick up the poles at $0,1,\infty$ and $\bar{z}$. It yields
\begin{equation}
\begin{aligned}
  0&=\frac{1}{b^2}\partial_z^2\,\mathcal{V}_{4}-\frac{\Delta_{2,1}}{z(z-1)}\mathcal{V}_{4}-\frac{2z-1}{z(z-1)}\partial_z \mathcal{V}_{4}-\frac{\Delta_1'}{z^2(z-1)}\mathcal{V}_{4}+\frac{\Delta_2'}{z(1-z)^2}\mathcal{V}_{4}-\frac{\Delta'_3}{z(1-z)}\mathcal{V}_{4}\\
  &-\braket{\Psi_{P_1}^{\sigma 1 \sigma 2}(0)\frac{1}{2\pi i}\int_{w \sim \bar{z}}\frac{w(w-1)}{z(z-1)}\frac{T(w)}{w-z} V_{\braket{2,1}}(z,\bar{z})\Psi_{P_2}^{\sigma 2 \sigma 3}(1)\Psi_{P_3}^{\sigma 3 \sigma 1}(\infty)}\,.
\end{aligned}
\end{equation}
As 
\begin{equation}
\begin{aligned}
  T(w)V_{\braket{2,1}}(z,\bar{z}) &\sim \frac{\Delta_{2,1}V_{\braket{2,1}}(z,\bar{z})}{(w-z)^2}+\frac{\partial_zV_{\braket{2,1}}(z,\bar{z})}{w-z}+\frac{\Delta_{2,1}V_{\braket{2,1}}(z,\bar{z})}{(w-\bar{z})^2}+\frac{\partial_{\bar{z}}V_{\braket{2,1}}(z,\bar{z})}{w-\bar{z}}\,,\\
  \frac{w(w-1)}{z(z-1)(w-z)}&=\frac{\bar{z}(\bar{z}-1)}{z(z-1)}\frac{1}{\bar{z}-z}+\frac{(2\bar{z}-1)(\bar{z}-z)-\bar{z}(\bar{z}-1)}{z(z-1)(\bar{z}-z)^2}\times (w-z)+o(w-z)\,,
\end{aligned}
\end{equation}
I finally obtain after picking up the pole at $\bar{z}$, 
\begin{equation}
\begin{aligned}
  0&=\Bigg[\frac{1}{b^2}\partial_z^2-\frac{\Delta_{2,1}}{z(z-1)}-\frac{2z-1}{z(z-1)}\partial_z -\frac{\Delta_1'}{z^2(z-1)}+\frac{\Delta_2'}{z(1-z)^2}-\frac{\Delta'_3}{z(1-z)}\\
  &+\frac{\bar{z}(\bar{z}-1)\partial_{\bar{z}}}{z(z-1)(z-\bar{z})}-\frac{(2\bar{z}-1)(\bar{z}-z)-\bar{z}(\bar{z}-1)}{z(z-1)(\bar{z}-z)^2}\Delta_{\braket{2,1}}\Bigg]\mathcal{V}_{4}
\end{aligned}
\end{equation}
which can be rewritten as
\begin{equation}
\begin{aligned}
  0&=\Bigg[\frac{1}{b^2}\partial_z^2-\frac{2\Delta_{2,1}}{z(z-1)}-\frac{2z-1}{z(z-1)}\partial_z -\frac{\Delta_1'}{z^2(z-1)}+\frac{\Delta_2'}{z(1-z)^2}-\frac{\Delta'_3}{z(1-z)}\\
  &+\frac{\bar{z}(\bar{z}-1)\partial_{\bar{z}}}{z(z-1)(z-\bar{z})}-\frac{\Delta_{2,1}}{(z-\bar{z})^2}\Bigg]\mathcal{V}_{4}
\end{aligned}
\end{equation}
As expected, it yields a second order PDE. This equation should also be obeyed by all the conformal blocks cited above.
\paragraph{Boundary degenate field}
The decoupling equation of the null vector associated to $\Psi^{\sigma_2 \sigma_3}_{\braket{3,1}}(z)$ can be written as~\cite{Belavin:1984vu}
\begin{equation}
  \braket{\Psi_{P_1}^{\sigma_1 \sigma_2}(z_1)\chi_{\braket{3,1}}^{\sigma_2 \sigma_3}(z)\Psi_{P_2}^{\sigma_3 \sigma_4}(z_2)\Psi_{P_3}^{\sigma_4 \sigma_1}(z_3)}=0\,,
\end{equation}
where
\begin{equation}
  \chi_{\braket{3,1}}^{\sigma_2 \sigma_3}(z) = (\Delta_{3,1}+2)L_{-3}^{(z)}\Psi^{\sigma_2 \sigma_3}_{\braket{3,1}}(z) -2 L^{(z)}_{-1}L^{(z)}_{-2}\Psi^{\sigma_2 \sigma_3}_{\braket{3,1}}(z)+\frac{1}{\Delta_{3,1}+1}\Psi^{\sigma_2 \sigma_3}_{\braket{3,1}}(z)\,,
\end{equation}
with $\Delta_{3,1}=-1-2b^2$. The decoupling of $\chi_{\braket{3,1}}^{\sigma_2 \sigma_3}(z)$ takes the form~\cite{Belavin:1984vu}
\begin{equation}
\begin{aligned}
  0&=\Bigg[\frac{1}{\Delta_{3,1}}\frac{\partial^3}{\partial z^3}-\sum_{i=1}^3 \frac{2\Delta_{3,1}\Delta_i}{(z-z_i)}-\sum_{i=1}^3\frac{\Delta_{3,1}}{(z-z_i)^2}\frac{\partial}{\partial z_i}\\
  &-\sum_{i=1}^3 \frac{2\Delta_i}{(z-z_i)^2}-\sum_{i=1}^3 \frac{2}{(z-z_i)}\frac{\partial^2}{\partial z \partial z_i}\Bigg]\braket{\Psi_{P_1}^{\sigma_1 \sigma_2}(z_1)\Psi^{\sigma_2 \sigma_3}_{\braket{3,1}}(z)\Psi_{P_2}^{\sigma_3 \sigma_4}(z_2)\Psi_{P_3}^{\sigma_4 \sigma_1}(z_3)}\,.
\end{aligned}
\end{equation}
We can eliminate $\partial_{z_i}$ using global ward identities and then set $z_1=0,\,z_2=1,\,z_3=\infty$ to obtain\PG{I believe the derivation of~\cite{Belavin:1984vu} is wrong, there are factors of two differences...}
\begin{equation}
\begin{aligned}
  0&=\Bigg[\frac{1}{\Delta_{3,1}}\frac{\partial^3}{\partial z^3}+2\left[\frac{1}{z}+\frac{1}{z-1}\right]\frac{\partial}{\partial z^2}+\sum_{i=1}^3 \frac{\Delta_{3,1}-2\Delta_i}{z^2}
  -\sum_{i=1}^3 \frac{2\Delta_{3,1}\Delta_i}{(z-z_i)}-\sum_{i=1}^3\frac{\Delta_{3,1}}{(z-z_i)^2}\frac{\partial}{\partial z_i}\\
  &-\sum_{i=1}^3 \frac{2\Delta_i}{(z-z_i)^2}-\sum_{i=1}^3 \frac{2}{(z-z_i)}\frac{\partial^2}{\partial z \partial z_i}\Bigg]\braket{\Psi_{P_1}^{\sigma_1 \sigma_2}(z_1)\Psi^{\sigma_2 \sigma_3}_{\braket{3,1}}(z)\Psi_{P_2}^{\sigma_3 \sigma_4}(z_2)\Psi_{P_3}^{\sigma_4 \sigma_1}(z_3)}\,.
\end{aligned}
\end{equation}

\appendix
\section{Derivations of the BPZ equation}
\paragraph{Notations and useful}
Let $\mathcal{Q}_i(z_i),\mathcal{Q}(z)$ a set of quasi-primary fields inserted at the points $z_i,z$. Additionnaly we introduce the useful relations 
\begin{equation}
  z_{\shortrightarrow i} :=z-z_i \qquad \Delta_{\rm TOT} = \Delta + \sum_{i=1}^{N} \Delta_i\,.
\end{equation}
The correlation function $\mathcal{V}_{N+1}:=\braket{\mathcal{Q}(z)\prod_{i=1}^{N}\mathcal{Q}_i(z_i)}$ satisfies the global conformal ward identities. The global ward identities can be written in a matrix form as
\begin{equation}
  \begin{pmatrix}
  1 & 1 & 1 & \dots \\
  z_{\shortrightarrow 1} & z_{\shortrightarrow 2} & z_{\shortrightarrow 3} & \dots \\
  z_{\shortrightarrow 1}^2 & z_{\shortrightarrow 2}^2 & z_{\shortrightarrow 3}^2 & \dots \\
  \end{pmatrix}
  \begin{pmatrix}
    \partial_{z_{\shortrightarrow 1}} \\
    \partial_{z_{\shortrightarrow 2}} \\
    \partial_{z_{\shortrightarrow 3}} \\
    \dots
  \end{pmatrix}
  \mathcal{V}_{N+1}
  = \begin{pmatrix}
    \partial_z \\
    -\Delta_{\rm TOT} \\
    -2\sum_{i=1}^N \Delta_i z_{\shortrightarrow i} \\
    \dots
  \end{pmatrix}
  \mathcal{V}_{N+1}
  \,.
\end{equation}
In the case $N=3$ we recognize the Vandermonde matrix which can be inverted as long as the $z_{\shortrightarrow i}$ are different. In such a case, we can invert the matrix to eliminate the $ \partial_{z_{\shortrightarrow i}}$ in terms of $\partial_z$. One obtains
\begin{equation}
  \partial_{z_{\shortrightarrow i}}=\frac{1}{\prod_{j\neq i} (z_j-z_i)}\left[\left(\prod_{j \neq i }z_{\shortrightarrow j}\right)\partial_z +\left(\sum_{j\neq i}z_{\shortrightarrow j}\right)\Delta_{\rm TOT}-2\sum_{j=1}^{3}\Delta_j z_{\shortrightarrow j}\right]\,.
\end{equation}
This holds when acting on $\mathcal{V}_{4}$ (i.e 4 point correlation function constituted of quasi-primaries). In deriving the BPZ equation of order $2,3$, it appears to be useful to rewrite the following operators in terms of $\partial_z$
\begin{equation}
  \mathcal{O}_1 :=-\sum_{i=1}^{3}\frac{\partial_{z_{\shortrightarrow i}}}{z_{\shortrightarrow i}}  \qquad \mathcal{O}_2  :=-\sum_{i=1}^{3}\frac{\partial_{z_{\shortrightarrow i}}}{z^2_{\shortrightarrow i}} \,.
\end{equation}
We find
\begin{equation}
  \begin{aligned}
    \mathcal{O}_1  &= -\left(\sum_{i=1}^{3}\frac{1}{z_{\shortrightarrow i}}\right)\partial_z + \sum_{i=1}^{3}\frac{2\Delta_i -\Delta_{\rm TOT}}{\prod_{j\neq i}z_{\shortrightarrow j}}\,,\\
    \mathcal{O}_2  &= -\left(\sum_{i=1}^3 \frac{1}{z_{\shortrightarrow i}^2} +\sum_{i=1}^3 \prod_{j \neq i}\frac{1}{z_{\shortrightarrow j}}\right)\partial_z-\frac{\Delta_{\rm TOT}\prod_{i=1}^{3}(\sum_{j \neq i} z_{\shortrightarrow j})- 2\left(\sum_{i=1}^{3}\Delta_i z_{\shortrightarrow i}\right)\sum_{i=1}^{3}(\prod_{j \neq i} z_{\shortrightarrow j})}{\prod_{i=1}^{3}z_{\shortrightarrow i}^2}\,.
  \end{aligned}
\end{equation}
In particular case of interest $z_{\shortrightarrow 1}=z,\,z_{\shortrightarrow 2}=z-1,\,z_{\shortrightarrow 3} = z-z_{\infty}$ with $z_{\infty}\to \infty$ and rewritting $\Delta_3 = \Delta_{\infty}$, this greatly simplifies leaving
\begin{equation}\label{eq:veryuseful}
  \begin{aligned}
    \mathcal{O}^{(\infty)}_1  &= -\left(\frac{1}{z}+\frac{1}{z-1}\right)\partial_z +\frac{2\Delta_{\infty} -\Delta_{\rm TOT}}{z(z-1)}\,,\\
    \mathcal{O}^{(\infty)}_2  &= -\left(\frac{1}{z^2}+\frac{1}{(z-1)^2}+\frac{1}{z(z-1)}\right)\partial_z+\frac{\left(2\Delta_{\infty}-\Delta_{\rm TOT}\right)(2z-1)}{z^2(z-1)^2}\,.
  \end{aligned}
\end{equation}
\paragraph{Second order BPZ equation}
The decoupling of the null vector associated to $V_{\braket{2,1}}(z)$ yields the following equation
for $\mathcal{V}_{4}(z,z_1,z_2,z_3)=\braket{V_{\braket{2,1}}(z)V_{\Delta_1}(z_1)V_{\Delta_2}(z_2)V_{\Delta_3}(z_3)}$
\begin{equation}
\begin{aligned}
  0 &=\frac{1}{b^2}\partial_z^2 \mathcal{V}_{4} +\braket{L_{-2}^{(z)}V_{\braket{2,1}}(z)V_{\Delta_1}(z_1)V_{\Delta_2}(z_2)V_{\Delta_3}(z_3)}\,\\
  &=\frac{1}{b^2}\partial_z^2 \mathcal{V}_{4} +\left[\sum_{i=1}^{3}\frac{\Delta_i}{z_{\shortrightarrow i}^2}+\frac{\partial_{z_i}}{z_{\shortrightarrow i}}\right]\mathcal{V}_{4}\,\\
  &=\frac{1}{b^2}\partial_z^2 \mathcal{V}_{4} +\left[\sum_{i=1}^{3}\frac{\Delta_i}{z_{\shortrightarrow i}^2}-\frac{\partial_{z_{\shortrightarrow i}}}{z_{\shortrightarrow i}}\right]\mathcal{V}_{4}\,\\
  &=\frac{1}{b^2}\partial_z^2 \mathcal{V}_{4} +\left[\sum_{i=1}^{3}\frac{\Delta_i}{z_{\shortrightarrow i}^2}+ \mathcal{O}_1\right]\mathcal{V}_{4}.\\
\end{aligned}
\end{equation}
In the second line, we deform the contour sitting at $z$ ($L_{-2}^{(z)}=\frac{1}{2\pi i}\int_{w \sim z}\frac{dw T(w)}{w-z}$) to pick up the poles at the other positions. In the third line, we used $\partial_{z_i}=-\partial_{z_{\shortrightarrow i}}$. Sending the insertion points to $0,1,\infty$ and defining the rescaled correlation function
\begin{equation}
  \mathcal{F}(z) := \underset{z_1 \to 0}{\lim} \,\underset{z_2 \to 1}{\lim}\, \underset{z_3 \to \infty}{\lim} \left(z_3^{2\Delta_{\infty}}\mathcal{V}_{4}(z,z_1,z_2,z_3)\right)\,,
\end{equation}
we obtain
\begin{equation}
  0 =\frac{1}{b^2}\partial_z^2 \mathcal{F}(z) +\left[\frac{\Delta_1}{z^2}+\frac{\Delta_2}{(z-1)^2}-\left(\frac{1}{z}+\frac{1}{z-1}\right)\partial_z +\frac{2\Delta_{\infty} -\Delta_{\rm TOT}}{z(z-1)}\right]\mathcal{F}(z).\\
\end{equation}
After massaging the last equation, this yields
\begin{equation}
  0 =\frac{z(1-z)}{b^2}\partial_z^2 \mathcal{F}(z) +(2z-1) \partial_z \mathcal{V}_{4} + \left[\Delta_{2,1}-\Delta_{\infty}+\frac{\Delta_1}{z}+\frac{\Delta_2}{1-z}\right]\mathcal{F}(z)\,.
\end{equation}
\paragraph{Third order BPZ equation}
The null vector is given by 
\begin{equation}\label{eq:decoupling3}
  \chi_{3,1}(z)=(\Delta_{3,1}+2)L_{-3} V_{3,1}(z) -2 L_{-1} L_{-2}V_{3,1}(z)+\frac{L_{-1}^3}{\Delta_{3,1}+1}V_{3,1}(z)\,.
\end{equation}
We already know from the previous paragraph that upon sending $z_1 \to 0, z_2 \to 1, z_{3}\to \infty$, the action of $L_{-2}$ reduces to
\begin{equation}\label{eq:l-2}
  \frac{\Delta_1}{z^2}+\frac{\Delta_2}{(z-1)^2}-\left(\frac{1}{z}+\frac{1}{z-1}\right)\partial_z +\frac{2\Delta_{\infty}-\Delta_{\rm TOT}}{z(z-1)}\,.
\end{equation}
It remains to know the action of $L_{-3}$. We have
\begin{equation}
  \begin{aligned}
    &\braket{L_{-3}^{(z)}V_{\braket{3,1}}(z)V_{\Delta_1}(z_1)V_{\Delta_2}(z_2)V_{\Delta_3}(z_3)}\\
    &=\braket{\frac{1}{2\pi i}\int _{w \sim z} \frac{T(w)}{(w-z)^2}V_{\braket{3,1}}(z)V_{\Delta_1}(z_1)V_{\Delta_2}(z_2)V_{\Delta_3}(z_3)}\\
    &=-\sum_{i=1}^3 \braket{V_{\braket{3,1}}(z) \frac{1}{2\pi i}\int _{w \sim z_i} \frac{T(w)}{(w-z)^2}V_{\Delta_i}(z_i) \dots}\\
    &=\sum_{i=1}^3 \braket{V_{\braket{3,1}}(z) \left[\frac{2\Delta_i}{(z_i-z)^3}-\frac{\partial_{z_i}}{(z_i-z)^2}\right]V_{\Delta_i}(z_i) \dots}\\
    &=-\sum_{i=1}^3 \braket{V_{\braket{3,1}}(z) \left[\frac{2\Delta_i}{z_{\shortrightarrow i}^3}-\frac{\partial_{z_{\shortrightarrow i}}}{z_{\shortrightarrow i}^2}\right]V_{\Delta_i}(z_i) \dots}\\
    &=-\left(\sum_{i=1}^3 \frac{2\Delta_i}{z_{\shortrightarrow i}^3}\mathcal{V}_{4}\right)-\mathcal{O}_{2} \mathcal{V}_{4}\,.
  \end{aligned}
\end{equation}
Sending $z_1 \to 0, z_2 \to 1, z_{3}\to \infty$ and using~\eqref{eq:veryuseful}, we obtain that the action of $L_{-3}$ on $\mathcal{F}(z)$ is
\begin{equation}\label{eq:l-3}
  -\frac{2\Delta_1}{z^3}-\frac{2\Delta_2}{(z-1)^3}+\left(\frac{1}{z^2}+\frac{1}{(z-1)^2}+\frac{1}{z(z-1)}\right)\partial_z -\frac{(2\Delta_{\infty}-\Delta_{\rm TOT})(2z-1)}{z^2(z-1)^2}\,.
\end{equation}
Combining~\eqref{eq:l-2}\,\eqref{eq:l-3} with~\eqref{eq:decoupling3}, we obtain
\begin{equation}
  \begin{aligned}
    0 &= \Bigg[\frac{\partial_z^3}{\Delta_{3,1}+1} + 2 \left[\frac{1}{z}+\frac{1}{z-1}\right]\partial^2_z \\
    +&\left[\frac{\Delta_{3,1}-2\Delta_1}{z^2} + \frac{\Delta_{3,1}-2\Delta_2}{(z-1)^2}+\frac{\Delta_{3,1}+2-2(2\Delta_{\infty}-\Delta_{\rm TOT})}{z(z-1)}\right]\partial_z \\
    -&\left[\frac{2\Delta_{3,1}\Delta_1}{z^3}+\frac{2\Delta_{3,1}\Delta_2}{(z-1)^3}+\frac{\Delta_{3,1}(2\Delta_{\infty}-\Delta_{\rm TOT})(2z-1)}{z^2(z-1)^2}\right]
    \Bigg] \mathcal{F}(z)\,.
  \end{aligned}
\end{equation}
The key check is the indicial equations at $0,1,\infty$ which should reproduce the results of the fusion rules. At $0,1$ the indicial equation is easily obtained by inserting a power ansatz $z^\alpha$/$(z-1)^\alpha$ and we easily find
\begin{equation}
  \alpha(\alpha-1)(\alpha-2)+2\alpha(\alpha-1)(\Delta_{3,1}+1)+\alpha(\Delta_{3,1}-2\Delta_i)(\Delta_{3,1}+1)-2\Delta_{3,1}\Delta_i(\Delta_{3,1}+1)=0\,,i\in\{1,2\}\,.
\end{equation}
It is then a matter of algebra to check that the solutions of this equation are
\begin{equation}
  \begin{aligned}
    \alpha_{0} &= \Delta_{i} - \Delta_{3,1} - \Delta_{i} = -\Delta_{3,1} \\
    \alpha_{i , \pm} &= \Delta_{i \pm b} - \Delta_{3,1} - \Delta_{i}\,,
  \end{aligned}
\end{equation}
$\Delta = \frac{(b+b^{-1})^2}{4}-P^2$ and $\pm$ means we shift the momentum by b. Recall $P_{r,s} = \frac{rb+sb^{-1}}{2}$. We then check the exponent at infinity. We first perform the substitution $u=\frac{1}{z}$. We have
\begin{equation}
\begin{aligned}
   \partial_z &= -u^2 \partial_u \\
   \partial_z^2 &= 2u^3 \partial_u +u^4 \partial_u^2 \\
   \partial_z^3 &= -6u^4 \partial_u -6u^5 \partial^2_u -u^6 \partial_u^3\,.\\
\end{aligned}
\end{equation}
After some algebra we obtain the equation
\begin{equation}
  \begin{aligned}
    0 &= \frac{\partial_u^3}{\Delta_{3,1}+1} + \frac{6}{u(\Delta_{3,1}+1)}\partial_u^2 +\frac{6}{u^2(\Delta_{3,1}+1)}\partial_u\\
    -&2\left[\frac{2}{u}+\frac{1}{1-u}\right]\left[\frac{2}{u}\partial_u +\partial_u^2\right] \\
    +&\left[\frac{\Delta_{3,1}-2\Delta_1}{u^2}+\frac{\Delta_{3,1}-2\Delta_2}{u^2(1-u)^2}+\frac{\Delta_{3,1}+2-2(2\Delta_{\infty}-\Delta_{\rm TOT})}{u^2(1-u)}\right]\partial_u\\
    +&\left[\frac{2\Delta_{3,1}\Delta_1}{u^3}+\frac{2\Delta_{3,1}\Delta_2}{u^3(1-u)^3}+\frac{\Delta_{3,1}(2\Delta_{\infty}-\Delta_{\rm TOT})(2-u)}{u^3(1-u)^2}\right]\,.
  \end{aligned}
\end{equation}
By looking to a power law ansatz $u^{\alpha}$, we find the indicial equation
\begin{equation}
  \alpha(\alpha+1)(\alpha+2)-4(\Delta_{3,1}+1)\alpha(\alpha+1)+\alpha(\Delta_{3,1}+1)(5\Delta_{3,1}+2 -2\Delta_{\infty})+2(\Delta_{3,1}+1)(\Delta_{\infty}-\Delta_{3,1}) =0\,.
\end{equation}
This equation should give the following exponents
\begin{equation}
  \begin{aligned}
    \alpha_{0} &= \Delta_{\infty} + \Delta_{3,1} - \Delta_{\infty} = \Delta_{3,1}\\
    \alpha_{+} &=  \Delta_{\infty +b} + \Delta_{3,1} - \Delta_{\infty} \\
    \alpha_{-} &=  \Delta_{\infty -b} + \Delta_{3,1} - \Delta_{\infty}\,.
  \end{aligned}
\end{equation}
Again, this is a simple matter of algebra that it is indeed the case.
\section{Boundary states: Ishibashi states vs Cardy states}
We first introduce Virasoro primary states. Highest-weight representations of the Virasoro algebra are generated from states $\ket{V_P}$ satisfying
\begin{equation}
  L_0 \ket{V_P} = \Delta_P \ket{V_P}, 
  \qquad 
  L_n \ket{V_P} = 0 \quad (n>0),
  \qquad 
  \Delta_P = \frac{Q^2}{4}-P^2 .
\end{equation}
We do not impose any normalization condition on the states $\ket{V_P}$.

On a conformal representative of a manifold with boundary, preserving one copy of the conformal symmetry requires the absence of momentum flow through the boundary,
\begin{equation}
  n_\mu T^{\mu\nu}=0 .
\end{equation}
At the quantum level, this condition translates into the gluing condition
\begin{equation}
  \left(L_n-\bar L_{-n}\right)\ket{B}=0,
  \qquad n\in \mathbb{Z},
\end{equation}
where the boundary state $\ket{B}$ is viewed as a closed-string state.

For each Virasoro representation generated by a primary state $\ket{V_P}$, there exists an Ishibashi state, denoted by $\ket{P}$, which solves the gluing condition. It can be written formally as
\begin{equation}
  \ket{P}
  =
  \sum_{N\geq 0}
  \sum_{I,J\in \mathcal{P}_N}
  \left(S^{(N)}\right)^{-1}_{IJ}
  L_{-I}\ket{V_P}
  \otimes
  \overline{L_{-J}\ket{V_P}} ,
\end{equation}
where $\mathcal{P}_N$ denotes the set of partitions of the level $N$, 
$L_{-I}=L_{-i_1}\cdots L_{-i_k}$ for $I=(i_1,\ldots,i_k)$ with $\sum_a i_a=N$, and $S^{(N)}$ is the Shapovalov matrix of Virasoro descendants at level $N$.
In general, physical boundary states also called Cardy states need not to be Ishibashi states. Physical boundary states $\ket{s_1},\,\ket{s_2}$ satisfy the open/close string duality Cardy condition. Let an annulus of unit width for which the circle has length $\beta$. Then
\begin{equation}
\begin{aligned}
  Z_{s_1\,s_2}(\beta) &= \rm{Tr}_{\mathcal{H}_{\rm open}}[e^{-\beta H_{s_1\,s_2}}] = \sum_{h} n^{s_1\,s_2}_h \chi_h(\beta) \\
                &= \bra{s_1} e^{-H/\beta} \ket{s_2}\,.
\end{aligned}
\end{equation}
This condition should be understood as determining $n_{h}^{s_1\,s_2}$.

\bibliographystyle{JHEP}
\bibliography{biblio}

\end{document}

